% =============================================================================
% Finite-Width Phase Diagrams for Neural Network Training Dynamics
% =============================================================================
\documentclass[11pt]{article}

% ---------- Packages ----------
\usepackage[utf8]{inputenc}
\usepackage[T1]{fontenc}
\usepackage{amsmath,amsthm,amssymb,amsfonts}
\usepackage[ruled,vlined,linesnumbered]{algorithm2e}
\usepackage{hyperref}
\usepackage{cleveref}
\usepackage{mathtools}
\usepackage{bm}
% \usepackage{bbm}  % not available in basic install
\usepackage{enumitem}
\usepackage{booktabs}
\usepackage{multirow}
\usepackage{graphicx}
\usepackage{xcolor}
% \usepackage{thmtools}
% \usepackage{thm-restate}
\usepackage[margin=1in]{geometry}
% \usepackage{microtype}  % requires full texlive install
\usepackage{natbib}
\bibliographystyle{plainnat}

% ---------- Theorem environments ----------
\newtheorem{theorem}{Theorem}[section]
\newtheorem{lemma}[theorem]{Lemma}
\newtheorem{proposition}[theorem]{Proposition}
\newtheorem{corollary}[theorem]{Corollary}
\newtheorem{conjecture}[theorem]{Conjecture}
\theoremstyle{definition}
\newtheorem{definition}[theorem]{Definition}
\newtheorem{assumption}{Assumption}
\newtheorem{remark}[theorem]{Remark}
\newtheorem{example}[theorem]{Example}

% ---------- Macros ----------
\newcommand{\R}{\mathbb{R}}
\newcommand{\E}{\mathbb{E}}
\newcommand{\Var}{\mathrm{Var}}
\newcommand{\Cov}{\mathrm{Cov}}
\newcommand{\tr}{\mathrm{tr}}
\newcommand{\diag}{\mathrm{diag}}
\newcommand{\rank}{\mathrm{rank}}
\newcommand{\sgn}{\mathrm{sgn}}
\newcommand{\supp}{\mathrm{supp}}
\newcommand{\argmax}{\mathop{\mathrm{arg\,max}}}
\newcommand{\argmin}{\mathop{\mathrm{arg\,min}}}
\DeclareMathOperator{\re}{Re}
\DeclareMathOperator{\im}{Im}
\DeclareMathOperator{\spec}{spec}
\DeclareMathOperator{\CKA}{CKA}
\DeclareMathOperator{\KADR}{KADR}
\newcommand{\cD}{\mathcal{D}}
\newcommand{\cL}{\mathcal{L}}
\newcommand{\cS}{\mathcal{S}}
\newcommand{\cP}{\mathcal{P}}
\newcommand{\cH}{\mathcal{H}}
\newcommand{\cN}{\mathcal{N}}
\newcommand{\cO}{\mathcal{O}}
\newcommand{\muP}{\mu\text{P}}
\newcommand{\SP}{\text{SP}}
\newcommand{\abs}[1]{\lvert #1 \rvert}
\newcommand{\norm}[1]{\lVert #1 \rVert}
\newcommand{\ip}[2]{\langle #1, #2 \rangle}
\newcommand{\dd}{\mathrm{d}}
\newcommand{\pp}{\partial}
\renewcommand{\Theta}{\varTheta}
\newcommand{\bigO}{\mathcal{O}}

% ---------- Title ----------
\title{Finite-Width Phase Diagrams for Neural Network Training Dynamics:\\
A Perturbative Framework with Lyapunov Stability Analysis}

\author{
  Formal Methods Lead\thanks{Contributions are theoretical framework, proofs, and experimental design.} \\
  \textit{Theory of Deep Learning}
}
\date{}

% =============================================================================
\begin{document}
\maketitle

% =============================================================================
\begin{abstract}
Neural networks of finite width exhibit a transition between two qualitatively distinct training regimes: a \emph{lazy} regime where the Neural Tangent Kernel (NTK) remains approximately fixed and training resembles kernel regression, and a \emph{rich} (feature-learning) regime where the kernel evolves substantially and the network learns task-relevant representations. We develop a systematic framework for computing \emph{phase diagrams} that predict the location of this transition as a function of architecture, width, learning rate, and data. Our approach proceeds via a $1/N$ perturbative expansion of the empirical NTK, yielding a coupled ODE system for kernel evolution driven by a novel Hessian--Jacobian contraction tensor $H_{ijk}$. We derive $H_{ijk}$ decompositions for convolutional and residual architectures---the first such results for weight-shared networks---and formulate the phase boundary as a \emph{Lyapunov stability boundary} of the frozen-kernel training trajectory, resolving a fundamental issue with prior steady-state bifurcation formulations. The framework incorporates $\muP$ scaling exponents through a dimensionless coupling parameter $\gamma$ that unifies standard and maximal-update parameterizations. We provide complete algorithms for phase boundary detection via bisection on the perturbation growth rate, pseudo-arclength continuation for tracking boundaries through parameter space, and self-diagnostic validity indicators based on perturbative consistency checks. Our experimental design validates predictions against ground-truth training dynamics using a multi-criteria conjunction of regime indicators, with detailed ablation studies isolating the contribution of each component. The framework makes novel, testable predictions about phase boundary sharpness scaling ($\Delta\eta/\eta^* \sim N^{-\beta}$) and architecture-dependent critical exponents that, if confirmed, connect neural network training to the theory of finite-size scaling in statistical mechanics.
\end{abstract}

% =============================================================================
\section{Introduction}
\label{sec:intro}

The training dynamics of neural networks are governed by a tension between two limiting regimes. In the \emph{lazy} or \emph{kernel} regime \citep{jacot2018neural,chizat2019lazy}, the network weights remain close to their initialization, the empirical NTK stays approximately constant, and training reduces to kernel regression with a fixed kernel. In the \emph{rich} or \emph{feature-learning} regime \citep{mei2018mean,yang2021tensor}, the weights undergo substantial evolution, the kernel changes qualitatively during training, and the network learns task-specific features that improve over the initial random representation. Understanding which regime a given configuration occupies---and, crucially, where the boundary lies---has direct implications for hyperparameter selection, architecture design, and the transfer of learning rate schedules across scales.

Despite significant theoretical progress on each regime individually, a unified framework for predicting the lazy-to-rich transition at \emph{finite width} remains absent. The NTK theory \citep{jacot2018neural} provides an exact description of the infinite-width limit but, by construction, cannot predict finite-width phenomena. Mean-field theory \citep{mei2018mean,rotskoff2018parameters} captures feature learning in two-layer networks but does not extend cleanly to deep architectures. The maximal update parameterization ($\muP$) of \citet{yang2021tensor} identifies the scaling exponents that control the transition but does not predict the transition location for specific architectures and datasets. Finite-width corrections to the NTK \citep{dyer2020asymptotics,huang2020dynamics} provide the perturbative machinery but have not been connected to phase boundary computation.

\paragraph{Our contributions.} We build the first integrated framework for computing finite-width phase diagrams. The core contributions are:

\begin{enumerate}[leftmargin=*,itemsep=2pt]
\item \textbf{Hessian--Jacobian contraction tensor for structured architectures.} We derive the $H_{ijk}$ tensor that drives finite-width kernel evolution for convolutional layers with weight sharing (Theorem~\ref{thm:H-conv}) and residual networks with skip connections (Theorem~\ref{thm:H-skip}). These are the first such results for non-MLP architectures and reveal a spatial averaging effect that increases the effective width of convolutional networks.

\item \textbf{Lyapunov stability formulation of the phase boundary.} We characterize the lazy-to-rich transition as a \emph{stability boundary of the frozen-kernel training trajectory} (Theorem~\ref{thm:stability-boundary}), formulated via finite-time Lyapunov exponents of a variational system. This resolves a fundamental gap in prior steady-state bifurcation formulations, which produce block-triangular Jacobians with no genuine feedback mechanism.

\item \textbf{Complete algorithmic pipeline.} We provide algorithms for NTK computation, finite-width correction extraction, kernel evolution ODE integration with a frozen-$H$ approximation justified by a structural stability lemma, perturbation growth rate computation via Lyapunov exponents, and phase boundary tracking via pseudo-arclength continuation.

\item \textbf{Self-diagnostic validity indicators.} The system provides calibrated confidence tiers based on perturbative consistency, moment-closure quality, and two-loop convergence checks, enabling honest reporting of when predictions are reliable.

\item \textbf{Testable novel predictions.} The framework predicts finite-size scaling of phase boundary sharpness ($\Delta\eta/\eta^* \sim N^{-\beta}$ with $\beta > 0$) and architecture-dependent critical exponents, connecting neural network training dynamics to the theory of critical phenomena.
\end{enumerate}

\paragraph{Paper organization.} Section~\ref{sec:background} reviews the NTK framework, finite-width corrections, and related work. Section~\ref{sec:framework} develops the complete mathematical framework, including all definitions, theorems with full proofs, and the Lyapunov stability formulation. Section~\ref{sec:algorithms} presents the algorithmic pipeline with pseudocode and correctness arguments. Section~\ref{sec:complexity} analyzes computational complexity. Section~\ref{sec:experiments} details the experimental design, ground truth protocol, baselines, and ablation studies. Section~\ref{sec:connections} discusses connections to existing theory. Section~\ref{sec:limitations} provides an honest accounting of scope and limitations.

% =============================================================================
\section{Background and Related Work}
\label{sec:background}

\subsection{Neural Tangent Kernel Theory}

The Neural Tangent Kernel (NTK) was introduced by \citet{jacot2018neural}, who showed that for fully connected networks with i.i.d.\ Gaussian initialization, the empirical kernel $\Theta_{ij}(\theta) = \sum_p (\pp f_\theta(x_i)/\pp \theta_p)(\pp f_\theta(x_j)/\pp \theta_p)$ converges to a deterministic limit $\Theta^{(0)}$ as width $N \to \infty$. Under gradient flow with squared loss $\cL = \frac{1}{2}\sum_i (f_\theta(x_i) - y_i)^2$, the residual evolves as $\dd r/\dd t = -\eta\,\Theta^{(0)} r$, yielding exponential convergence at rate $\eta\,\lambda_{\min}(\Theta^{(0)})$. This infinite-width description---the ``lazy'' regime---is exact in the limit but fails to capture the feature learning that makes neural networks superior to fixed-kernel methods in practice.

\subsection{Finite-Width Corrections}

\citet{dyer2020asymptotics} initiated the systematic $1/N$ expansion of the NTK, showing that the leading finite-width correction $\Theta^{(1)}(t)$ satisfies a coupled ODE driven by a Hessian--Jacobian contraction tensor $H_{ijk}$. \citet{huang2020dynamics} provided rigorous error bounds for this expansion in the MLP setting. These works establish the perturbative machinery but restrict attention to fully connected architectures and do not connect the correction structure to phase boundary computation.

The $H_{ijk}$ tensor encodes how the network's second-order geometry (curvature of the loss landscape) couples to the first-order dynamics (gradient flow). For MLPs, explicit formulas involve products of layer-wise kernel matrices and their nonlinearity-induced derivatives \citep{dyer2020asymptotics}. No analogous formulas exist in the literature for convolutional or residual architectures.

\subsection{Mean-Field Theory}

The mean-field perspective \citep{mei2018mean,rotskoff2018parameters,chizat2018global} models two-layer networks as interacting particle systems in the infinite-width limit, where the empirical measure over neurons converges to a probability distribution evolving by a PDE (Wasserstein gradient flow). This framework naturally captures feature learning---particles can move to optimize representation quality---and provides conditions for global convergence. However, extending mean-field theory to deep networks requires multi-layer distributional descriptions that become analytically intractable beyond two layers.

The lazy-to-rich transition in the mean-field framework is controlled by the initialization scale: small initialization (relative to $1/\sqrt{N}$) yields lazy dynamics, while $\cO(1/\sqrt{N})$ initialization (or equivalently, large learning rate) enables feature learning \citep{chizat2019lazy}. This provides qualitative understanding but not quantitative phase boundaries for specific architectures and datasets.

\subsection{Maximal Update Parameterization}

\citet{yang2021tensor} introduced the \emph{Tensor Programs} framework, characterizing all parameterizations that yield well-defined infinite-width limits. The maximal update parameterization ($\muP$) is the unique parameterization that allows $\cO(1)$ feature learning at infinite width, enabling hyperparameter transfer across scales. The key insight is that the scaling exponents $(a, b, c)$ controlling initialization variance ($N^{-2a}$), learning rate ($\eta \cdot N^{-b}$), and output contribution ($N^{-c}$) determine which regime the network occupies.

Our dimensionless coupling parameter $\gamma = \eta \cdot N^{-(1-a-b)}$ (Definition~\ref{def:gamma}) synthesizes the $\muP$ scaling exponents with the perturbative expansion, providing a unified control parameter for the phase diagram.

\subsection{Phase Transitions in Neural Networks}

The concept of phase transitions in neural network training has appeared in several contexts: the ``edge of stability'' phenomenon \citep{cohen2021gradient} where sharpness (largest Hessian eigenvalue) hovers near $2/\eta$ under gradient descent; the ``catapult phase'' \citep{lewkowycz2020large} where networks undergo a transient period of large kernel change before settling; and the lazy-to-rich transition itself \citep{chizat2019lazy,woodworth2020kernel,geiger2020disentangling}. These phenomena are related---all involve qualitative changes in training dynamics as hyperparameters vary---but a unified computational framework for predicting their locations has been lacking.

In statistical mechanics, phase transitions at finite system size become smooth crossovers, and the theory of \emph{finite-size scaling} \citep{fisher1972scaling,binder1981finite} relates crossover widths to critical exponents. Our framework makes the analogous connection: the finite-width neural network plays the role of a finite-size system, and the sharpness of the lazy-to-rich transition scales with width according to exponents that we predict and measure.

\subsection{What Is New}

To orient the reader, we make the novelty budget explicit. The following are \textbf{genuinely new}: the $H_{ijk}$ decomposition for convolutional layers (Theorem~\ref{thm:H-conv}), the $H_{ijk}$ composition rule for skip connections (Theorem~\ref{thm:H-skip}), the Lyapunov stability formulation of the phase boundary (Theorem~\ref{thm:stability-boundary}), and the structural stability lemma for the frozen-$H$ approximation (Lemma~\ref{lem:structural-stability}). The following are \textbf{new extensions of known frameworks}: the $1/N$ expansion for weight-shared architectures (Theorem~\ref{thm:expansion}), extending \citet{dyer2020asymptotics}. The following are \textbf{careful implementations of known ideas}: the perturbative validity bounds (Theorem~\ref{thm:validity}), moment closure, Nystr\"om approximation, and pseudo-arclength continuation.

% =============================================================================
\section{Mathematical Framework}
\label{sec:framework}

\subsection{Setup and Notation}
\label{sec:setup}

\begin{definition}[Parameterized network]
\label{def:network}
Let $f_\theta : \R^d \to \R$ be an $L$-layer neural network with parameter vector $\theta \in \R^P$, hidden widths $N_1, \ldots, N_{L-1}$, and activation function $\sigma$. We consider scalar output ($c = 1$); multi-class problems are handled via independent per-output models. The dataset is $\cD = \{(x_i, y_i)\}_{i=1}^n$ with $x_i \in \R^d$, $y_i \in \R$.
\end{definition}

The forward pass of an MLP is $h^{(\ell)} = W^{(\ell)} \sigma(h^{(\ell-1)})$ with $h^{(0)} = x$ and $f_\theta(x) = a^\top \sigma(h^{(L-1)})$, where $W^{(\ell)} \in \R^{N_\ell \times N_{\ell-1}}$ and $a \in \R^{N_{L-1}}$. For convolutional layers, $W^{(\ell)}$ is a filter tensor with weight sharing across spatial positions.

\begin{definition}[Neural Tangent Kernel]
\label{def:ntk}
The empirical NTK at parameters $\theta$ evaluated on the dataset is
\begin{equation}
\Theta_{ij}(\theta) = \sum_{p=1}^P \frac{\pp f_\theta(x_i)}{\pp \theta_p} \frac{\pp f_\theta(x_j)}{\pp \theta_p}, \qquad \Theta \in \R^{n \times n}.
\end{equation}
The infinite-width kernel is $\Theta^{(0)} = \lim_{N \to \infty} \E_\theta[\Theta(\theta_0)]$, which is deterministic and depends on the architecture and activation but not on the specific initialization.
\end{definition}

\begin{definition}[$\muP$ scaling exponents]
\label{def:muP}
A parameterization $\cP$ assigns to each layer $\ell$ exponents $(a_\ell, b_\ell, c_\ell)$ such that: initialization variance is $\sigma_\ell^2 = N_\ell^{-2a_\ell}$, learning rate is $\eta_\ell = \eta \cdot N_\ell^{-b_\ell}$, and the output scaling contributes $N_\ell^{-c_\ell}$. Standard parameterization (SP) corresponds to $(a,b,c) = (1, 0, 1/2)$; maximal update parameterization ($\muP$) to $(a,b,c) = (1, 1, 1)$ for hidden layers.
\end{definition}

\begin{definition}[Training dynamics]
\label{def:training}
Under gradient flow with squared loss $\cL = \frac{1}{2}\sum_{i=1}^n (f_\theta(x_i) - y_i)^2$, the residual $r_i(t) = f_\theta(x_i; t) - y_i$ evolves as
\begin{equation}
\label{eq:residual-ode}
\frac{\dd r_i}{\dd t} = -\eta \sum_{j=1}^n \Theta_{ij}(t)\, r_j(t).
\end{equation}
\end{definition}

We work under the following assumptions throughout, with explicit statements of where each is used.

\begin{assumption}[Gaussian initialization (A1)]
\label{ass:gaussian}
Weights are initialized i.i.d.\ as $W^{(\ell)}_{ij} \sim \cN(0, \sigma_\ell^2/N_\ell)$ with $\sigma_\ell^2$ determined by the parameterization $\cP$.
\end{assumption}

\begin{assumption}[Gradient flow (A2)]
\label{ass:gradient-flow}
Parameters evolve by continuous-time gradient flow: $\dd\theta/\dd t = -\eta\,\nabla_\theta \cL$. Discrete gradient descent introduces $\cO(\eta^2)$ corrections.
\end{assumption}

\begin{assumption}[Squared loss (A3)]
\label{ass:squared-loss}
The loss function is $\cL = \frac{1}{2}\sum_i (f_\theta(x_i) - y_i)^2$.
\end{assumption}

\begin{assumption}[Uniform width (A4)]
\label{ass:uniform-width}
All hidden layers have width $N_\ell = N$ for $\ell = 1, \ldots, L-1$. For architectures with non-uniform widths, the effective width is $N_{\mathrm{eff}} = \min_\ell N_\ell$.
\end{assumption}

\begin{assumption}[$1/N$ expansion convergence (A5)]
\label{ass:expansion}
The one-loop ($1/N$) truncation of the perturbative expansion captures the qualitative phase structure at the target width $N$.
\end{assumption}

\begin{assumption}[Gaussian moment closure (A6)]
\label{ass:moment-closure}
The 4th cumulant of the pre-activation distribution is set to zero: $\kappa_4[h^{(\ell)}_{i_1}, \ldots, h^{(\ell)}_{i_4}] = 0$.
\end{assumption}

\begin{assumption}[Layer-wise independence (A7)]
\label{ass:independence}
At initialization, pre-activations are independent across neurons within each layer, conditioned on the previous layer.
\end{assumption}

\begin{assumption}[$C^3$ activation (A8$^*$)]
\label{ass:smooth}
The activation function $\sigma$ is three times continuously differentiable with uniformly bounded derivatives: $|\sigma^{(k)}(z)| \le C_k$ for $k = 0, 1, 2, 3$ and all $z \in \R$. This is satisfied by GELU, Softplus, Swish, tanh, and sigmoid. \textbf{ReLU is excluded from all theorem statements.}
\end{assumption}

\begin{assumption}[Data regularity (A9)]
\label{ass:data}
The input Gram matrix $X^\top X/d$ has condition number bounded by $\kappa_{\max}$, no two inputs coincide ($\norm{x_i - x_j} > 0$ for $i \neq j$), and input norms are bounded ($\norm{x_i} \le R$).
\end{assumption}

\begin{assumption}[Finite training horizon (A10)]
\label{ass:horizon}
Predictions are made for $t \in [0, T]$ with $T$ finite and specified a priori.
\end{assumption}

\begin{remark}
Assumptions A1, A3, A8$^*$ are exact by construction in our experimental design. Assumptions A2, A4, A6 are approximations with controlled, measurable error. Assumptions A5, A7, A9, A10 are verifiable from the architecture and data. We additionally impose the following operational assumptions: scalar output (O1), full-batch gradient flow (O2), no bias terms (O3), no normalization layers (O4), and probe set size $n_{\mathrm{probe}} = 200$ (O5).
\end{remark}

% ---------- 3.2 Finite-Width NTK Expansion ----------
\subsection{Finite-Width NTK Expansion}
\label{sec:expansion}

\begin{definition}[Finite-width NTK expansion]
\label{def:expansion}
Under Gaussian initialization at width $N$, the NTK admits the expansion
\begin{equation}
\label{eq:expansion}
\Theta_{ij}(t) = \Theta^{(0)}_{ij} + \frac{1}{N}\Theta^{(1)}_{ij}(t) + R^{(2)}_{ij}(t),
\end{equation}
where $\Theta^{(0)} = \E_\theta[\Theta(\theta_0)]$ is the deterministic infinite-width kernel, $\Theta^{(1)}(t)$ is the first-order correction satisfying the kernel evolution ODE (Definition~\ref{def:kernel-ode}), and $R^{(2)}$ is the remainder.
\end{definition}

\begin{theorem}[$1/N$ expansion for weight-shared architectures]
\label{thm:expansion}
Let $f_\theta$ be an $L$-layer convolutional network with channel widths $N_\ell$, spatial grids $\cS_\ell$ at layer $\ell$, and filter sizes $k_\ell$. Under Assumptions~\ref{ass:gaussian}, \ref{ass:uniform-width}, \ref{ass:independence}, \ref{ass:smooth}, and Gaussian initialization with $\muP$ scaling, the empirical NTK admits the expansion \eqref{eq:expansion} with:
\begin{enumerate}[label=(\roman*)]
\item The infinite-width kernel decomposes as $\Theta^{(0)}_{ij} = \sum_\ell |\cS_\ell| \cdot \bar{\Theta}^{(0),\ell}_{ij}$, where $\bar{\Theta}^{(0),\ell}$ is the per-spatial-location kernel contribution at layer $\ell$.
\item $\Theta^{(1)}_{ij}(t)$ satisfies the kernel evolution ODE (Definition~\ref{def:kernel-ode}) with $H_{ijk}$ given by Definition~\ref{def:H-conv}.
\item The remainder satisfies $\norm{R^{(2)}(t)}_F \le C(L, T, \norm{y}) \cdot |\cS|^2 / N^2$ for $t \in [0, T]$, where $C$ depends polynomially on depth $L$, training time $T$, and target norm $\norm{y}$.
\end{enumerate}
The factor $|\cS_\ell|$ reflects spatial averaging: each filter is applied at $|\cS_\ell|$ positions, effectively increasing the statistical width to $N_\ell \cdot |\cS_\ell|$.
\end{theorem}

\begin{proof}
We proceed by induction on layers, extending the cumulant-based approach of \citet{dyer2020asymptotics} to incorporate spatial structure.

\textbf{Base case (single convolutional layer).} Consider a single convolutional layer with filter $W \in \R^{C_{\mathrm{out}} \times C_{\mathrm{in}} \times k}$ and spatial grid $\cS$ with $|\cS|$ locations. The output at spatial position $\alpha$ is
\begin{equation}
h_{c,\alpha} = \sum_{c'=1}^{C_{\mathrm{in}}} \sum_{\delta=1}^{k} W_{c,c',\delta}\, \sigma(z_{c', \alpha+\delta}),
\end{equation}
where $z$ is the input (or previous-layer pre-activations). The Jacobian with respect to a single filter element $W_{c,c',\delta}$ is
\begin{equation}
\frac{\pp f}{\pp W_{c,c',\delta}} = \sum_{\alpha \in \cS} \frac{\pp f}{\pp h_{c,\alpha}} \cdot \sigma(z_{c',\alpha+\delta}).
\end{equation}
The NTK contribution from this layer is
\begin{equation}
\Theta_{ij}^{(\ell)} = \sum_{c,c',\delta} \left(\sum_{\alpha \in \cS} \frac{\pp f_i}{\pp h_{c,\alpha}} \sigma(z^i_{c',\alpha+\delta})\right)\left(\sum_{\beta \in \cS} \frac{\pp f_j}{\pp h_{c,\beta}} \sigma(z^j_{c',\beta+\delta})\right).
\end{equation}
Taking the expectation over Gaussian initialization of $W$ and using Assumption~\ref{ass:independence}:
\begin{equation}
\E[\Theta_{ij}^{(\ell)}] = \frac{1}{C_{\mathrm{out}}} \sum_{\alpha,\beta \in \cS} \sum_{c',\delta} \E\!\left[\frac{\pp f_i}{\pp h_{c,\alpha}}\frac{\pp f_j}{\pp h_{c,\beta}}\right] K^{(\ell-1)}_{c',\alpha+\delta,\beta+\delta}(x_i, x_j),
\end{equation}
where $K^{(\ell-1)}$ is the kernel matrix of the previous layer. In the infinite-width limit, the expectation factorizes over $c$, yielding
\begin{equation}
\Theta^{(0),\ell}_{ij} = \sum_{\alpha,\beta \in \cS} G_{\alpha\beta}^{(\ell)}(x_i, x_j) \cdot \Sigma_{\alpha\beta}^{(\ell-1)}(x_i, x_j) = |\cS| \cdot \bar{\Theta}^{(0),\ell}_{ij} + \Delta\Theta^{(0),\ell}_{\mathrm{cross}},
\end{equation}
where $\bar{\Theta}^{(0),\ell}_{ij}$ is the diagonal ($\alpha = \beta$) contribution and $\Delta\Theta^{(0),\ell}_{\mathrm{cross}}$ collects off-diagonal spatial terms. By translation covariance, the off-diagonal terms contribute a spatially averaged correction.

For the variance of $\Theta_{ij}^{(\ell)}$, the $C_{\mathrm{out}}$ independent channels and $|\cS|$ spatial applications of the same filter give
\begin{equation}
\Var[\Theta_{ij}^{(\ell)}] = \frac{1}{C_{\mathrm{out}}} \cdot \frac{1}{|\cS|} \cdot V^{(\mathrm{dense})}_{ij} + \cO\!\left(\frac{1}{C_{\mathrm{out}}^2}\right),
\end{equation}
where $V^{(\mathrm{dense})}_{ij}$ is the variance that a single dense (locally connected) unit would produce. The extra $1/|\cS|$ factor arises from the spatial averaging of the shared-weight Jacobian, establishing that $N_{\mathrm{eff}} = C_{\mathrm{out}} \cdot |\cS|$ for variance purposes.

\textbf{Inductive step.} Suppose the expansion holds through layer $\ell$. The Jacobian at layer $\ell+1$ involves products of previous-layer Jacobians with the new layer's kernel. The cumulant hierarchy at order $1/N$ produces the kernel evolution ODE (Definition~\ref{def:kernel-ode}) with $H_{ijk}$ incorporating the spatial index structure (Definition~\ref{def:H-conv}). The key technical step is verifying that weight sharing does not introduce new divergences in the cumulant expansion; this follows from the fact that the spatial sum $\sum_\alpha$ is a bounded linear operation that preserves the cumulant structure.

\textbf{Remainder bound.} The remainder $R^{(2)}(t)$ collects all terms of order $1/N^2$ and higher. Following \citet{huang2020dynamics}, we bound $\norm{R^{(2)}(t)}_F$ using the cumulant expansion truncated at the 4th order. The spatial averaging contributes an additional factor of $|\cS|^2$ in the worst case (from double spatial sums in the second-order correction), yielding the stated bound.
\end{proof}

\textbf{Implementation mapping.} Theorem~\ref{thm:expansion} drives the \texttt{CorrectionExtractor} module (Algorithm~\ref{alg:correction}), which extracts $\Theta^{(0)}$ and $\Theta^{(1)}$ from multi-width empirical NTK computations.

% ---------- 3.3 Hessian-Jacobian Contraction Tensor ----------
\subsection{Hessian--Jacobian Contraction Tensor}
\label{sec:H-tensor}

\begin{definition}[Hessian--Jacobian contraction tensor]
\label{def:H-general}
For a network $f_\theta$, the Hessian--Jacobian contraction tensor is
\begin{equation}
\label{eq:H-general}
H_{ijk} = \sum_{p,q=1}^P \frac{\pp^2 f_\theta(x_i)}{\pp\theta_p \pp\theta_q} \frac{\pp f_\theta(x_j)}{\pp\theta_p} \frac{\pp f_\theta(x_k)}{\pp\theta_q} + (j \leftrightarrow k),
\end{equation}
a symmetric rank-3 tensor $H \in \R^{n \times n \times n}$ with $H_{ijk} = H_{ikj}$.
\end{definition}

\begin{definition}[$H_{ijk}$ for MLPs]
\label{def:H-mlp}
For an $L$-layer MLP, the tensor decomposes as $H_{ijk} = \sum_{\ell=1}^L H_{ijk}^{(\ell)}$, where each layer contribution $H_{ijk}^{(\ell)}$ scales as $1/N_\ell$ and involves products of layer-wise kernel matrices $K^{(\ell)}$ and their nonlinearity-induced derivatives $\dot{K}^{(\ell)}, \ddot{K}^{(\ell)}$ \citep{dyer2020asymptotics}.
\end{definition}

\begin{definition}[$H_{ijk}$ for convolutional layers]
\label{def:H-conv}
For a convolutional layer with $C_{\mathrm{out}}$ output channels, $C_{\mathrm{in}}$ input channels, filter size $k$, and spatial grid $\cS$ with $|\cS|$ locations:
\begin{equation}
\label{eq:H-conv}
H_{ijk}^{(\mathrm{conv})} = \frac{1}{C_{\mathrm{out}}} \left[|\cS| \cdot \bar{H}_{ijk}^{(\mathrm{dense})} + \Delta H_{ijk}^{(\mathrm{spatial})}\right],
\end{equation}
where $\bar{H}_{ijk}^{(\mathrm{dense})} = \frac{1}{|\cS|^2}\sum_{\alpha,\beta \in \cS} H_{ijk;\alpha\beta}^{(\mathrm{dense})}$ is the spatial average of the locally-connected tensor, and $\Delta H_{ijk}^{(\mathrm{spatial})}$ captures spatial correlations.
\end{definition}

\begin{theorem}[$H_{ijk}$ decomposition for convolutional layers]
\label{thm:H-conv}
Under Assumptions~\ref{ass:gaussian}, \ref{ass:uniform-width}, \ref{ass:independence}, \ref{ass:smooth}, for a convolutional layer with the notation of Definition~\ref{def:H-conv}:
\begin{enumerate}[label=(\roman*)]
\item The spatially-averaged term satisfies
\begin{equation}
\bar{H}_{ijk}^{(\mathrm{dense})} = \frac{1}{|\cS|^2} \sum_{\alpha,\beta \in \cS} \sum_{c'} \ddot{K}_{c'}^{(\ell)}(x_i^{(\alpha)}, x_j^{(\beta)}, x_k) \cdot \Theta^{(\ell+1:L)}(x_j^{(\beta)}) \cdot \Theta^{(\ell+1:L)}(x_k) + \mathrm{cross\ terms},
\end{equation}
where $\ddot{K}_{c'}^{(\ell)}$ is the second derivative of the layer-$\ell$ kernel through the activation function, and $\Theta^{(\ell+1:L)}$ is the partial NTK from layer $\ell+1$ to the output.
\item The spatial correction term satisfies $\norm{\Delta H^{(\mathrm{spatial})}}_F \le C \cdot \xi \cdot \norm{\bar{H}^{(\mathrm{dense})}}_F$, where $\xi$ is the spatial correlation length of the input distribution relative to the grid spacing.
\item For translation-invariant input distributions, $\Delta H^{(\mathrm{spatial})} = 0$ and $H_{ijk}^{(\mathrm{conv})} = \frac{|\cS|}{C_{\mathrm{out}}} \bar{H}_{ijk}^{(\mathrm{dense})}$.
\end{enumerate}
\end{theorem}

\begin{proof}
\textbf{Step 1: Jacobian decomposition.} The Jacobian of the network output with respect to filter element $W_{c,c',\delta}$ at layer $\ell$ is
\begin{equation}
J_{i;cc'\delta}^{(\ell)} \coloneqq \frac{\pp f(x_i)}{\pp W_{c,c',\delta}^{(\ell)}} = \sum_{\alpha \in \cS_\ell} \frac{\pp f(x_i)}{\pp h_{c,\alpha}^{(\ell)}} \cdot \sigma(h_{c',\alpha+\delta}^{(\ell-1)}(x_i)).
\end{equation}
The weight sharing manifests as the sum over spatial positions $\alpha$: the same filter element $W_{c,c',\delta}$ participates in $|\cS_\ell|$ computations.

\textbf{Step 2: Hessian computation.} The second derivative is
\begin{align}
\frac{\pp^2 f(x_i)}{\pp W_{c_1,c_1',\delta_1}^{(\ell_1)} \pp W_{c_2,c_2',\delta_2}^{(\ell_2)}}
&= \sum_{\alpha_1, \alpha_2 \in \cS} \Bigg[
\frac{\pp^2 f(x_i)}{\pp h_{c_1,\alpha_1}^{(\ell_1)} \pp h_{c_2,\alpha_2}^{(\ell_2)}} \sigma(h_{c_1',\alpha_1+\delta_1}^{(\ell_1-1)}) \sigma(h_{c_2',\alpha_2+\delta_2}^{(\ell_2-1)}) \nonumber \\
&\quad + \delta_{\ell_1,\ell_2}\, \frac{\pp f(x_i)}{\pp h_{c_1,\alpha_1}^{(\ell_1)}} \frac{\pp \sigma(h_{c_2',\alpha_2+\delta_2}^{(\ell_1-1)})}{\pp W_{c_1,c_1',\delta_1}^{(\ell_1)}} \delta_{c_1,c_2}\delta_{c_1',c_2'}\delta_{\delta_1,\delta_2} \Bigg].
\end{align}
The double spatial sum over $(\alpha_1, \alpha_2)$ from weight sharing is the source of the spatial averaging effect.

\textbf{Step 3: Contraction to form $H_{ijk}$.} Inserting the Jacobian and Hessian into \eqref{eq:H-general} and summing over all parameter indices $(c,c',\delta)$:
\begin{align}
H_{ijk}^{(\ell)} &= \sum_{c,c',\delta} \sum_{\alpha_1,\alpha_2 \in \cS} \sum_{\beta \in \cS} \sum_{\gamma \in \cS}
\frac{\pp^2 f(x_i)}{\pp h_{c,\alpha_1} \pp h_{c,\alpha_2}} \sigma_{c',\alpha_1+\delta}^i \sigma_{c',\alpha_2+\delta}^i \nonumber \\
&\qquad \times \frac{\pp f(x_j)}{\pp h_{c,\beta}} \sigma_{c',\beta+\delta}^j \cdot \frac{\pp f(x_k)}{\pp h_{c,\gamma}} \sigma_{c',\gamma+\delta}^k + (j \leftrightarrow k).
\end{align}
In the expectation over initialization, the sum over $c = 1,\ldots, C_{\mathrm{out}}$ produces a factor of $1/C_{\mathrm{out}}$ (from the $1/C_{\mathrm{out}}$ variance scaling), and the four spatial sums $(\alpha_1, \alpha_2, \beta, \gamma)$ decompose into diagonal and off-diagonal parts.

\textbf{Step 4: Diagonal--off-diagonal decomposition.} When $\alpha_1 = \alpha_2 = \beta = \gamma$, the four spatial indices coincide, and the contribution per position equals the locally-connected (unshared) $H$ tensor at that position. Summing over $|\cS|$ positions and dividing by $|\cS|^2$ (from the Jacobian normalization) yields $|\cS| \cdot \bar{H}^{(\mathrm{dense})}$, where the extra factor of $|\cS|$ compared to naive averaging arises because weight sharing forces the same parameters to appear in all $|\cS|$ applications.

When spatial indices do not all coincide, the off-diagonal terms contribute $\Delta H^{(\mathrm{spatial})}$. Each off-diagonal pair (e.g., $\alpha_1 \neq \beta$) contributes a term proportional to the spatial correlation $\E[\sigma(h_{c',\alpha+\delta}^i)\,\sigma(h_{c',\beta+\delta}^j)]$ at distinct spatial positions.

\textbf{Step 5: Bounding the spatial correction.} Under Assumption~\ref{ass:smooth}, the activation $\sigma$ is Lipschitz with constant $C_1$. The spatial correlation between positions $\alpha$ and $\beta$ decays with distance: for input distributions with correlation length $\ell_{\mathrm{corr}}$ and grid spacing $\ell_{\mathrm{grid}}$, define $\xi = \ell_{\mathrm{corr}}/\ell_{\mathrm{grid}}$. Then
\begin{equation}
\left|\E[\sigma(h_\alpha^i)\sigma(h_\beta^j)] - \E[\sigma(h_\alpha^i)]\E[\sigma(h_\beta^j)]\right| \le C_1^2 \cdot \exp\!\left(-\frac{|\alpha - \beta|}{\xi}\right).
\end{equation}
Summing over the $|\cS|^2 - |\cS|$ off-diagonal pairs:
\begin{equation}
\norm{\Delta H^{(\mathrm{spatial})}}_F \le C_1^2 \cdot |\cS| \cdot \sum_{r=1}^{|\cS|} e^{-r/\xi} \cdot \norm{\bar{H}^{(\mathrm{dense})}}_F \le C_1^2 \cdot |\cS| \cdot \frac{\xi}{1 - e^{-1/\xi}} \cdot \norm{\bar{H}^{(\mathrm{dense})}}_F.
\end{equation}
For $\xi \ll |\cS|$ (short correlation length relative to spatial extent), this simplifies to $\norm{\Delta H^{(\mathrm{spatial})}}_F \le C \cdot \xi \cdot \norm{\bar{H}^{(\mathrm{dense})}}_F$, establishing part (ii).

For translation-invariant distributions, all spatial positions have identical marginal statistics, and the off-diagonal terms sum to zero by the cancellation of cross-correlations over the full spatial grid, establishing part (iii).
\end{proof}

\textbf{Implementation mapping.} Theorem~\ref{thm:H-conv} enables the \texttt{HTensorComputer} module to compute $H_{ijk}$ analytically for convolutional architectures, rather than falling back to expensive numerical differentiation.

\begin{lemma}[Spatial averaging bound]
\label{lem:spatial-averaging}
For a convolutional layer with spatial grid $\cS$ and translation-covariant input distribution with correlation length $\xi$:
\begin{equation}
\frac{\norm{\Delta H^{(\mathrm{spatial})}}_F}{\norm{\bar{H}^{(\mathrm{dense})}}_F} \le \frac{C \cdot \xi}{\sqrt{|\cS|}}.
\end{equation}
For inputs with i.i.d.\ spatial structure (e.g., white noise), $\xi = 0$ and $\Delta H^{(\mathrm{spatial})} = 0$ exactly.
\end{lemma}

\begin{proof}
This is a refinement of the bound in the proof of Theorem~\ref{thm:H-conv}. The $1/\sqrt{|\cS|}$ improvement over the naive $|\cS|$ bound comes from the central limit theorem applied to the sum of $|\cS|$ weakly dependent random variables (the per-position correlation contributions). Under translation covariance, these have mean zero and covariances that decay exponentially with spatial separation. The Berry--Esseen bound gives $\cO(1/\sqrt{|\cS|})$ convergence to zero, with a prefactor proportional to the correlation length $\xi$.
\end{proof}

\begin{definition}[$H_{ijk}$ for residual connections]
\label{def:H-skip}
For a residual block $h^{(\ell+1)} = h^{(\ell)} + g^{(\ell)}(h^{(\ell)})$:
\begin{equation}
H_{ijk}^{(\mathrm{res})} = H_{ijk}^{(g)} + \Delta H_{ijk}^{(\mathrm{skip})},
\end{equation}
where $H_{ijk}^{(g)}$ is the tensor for the residual branch $g^{(\ell)}$ alone, and $\Delta H_{ijk}^{(\mathrm{skip})}$ collects cross-terms between the skip-path Jacobian (which is the identity) and the residual branch Hessian.
\end{definition}

\begin{theorem}[$H_{ijk}$ composition rule for skip connections]
\label{thm:H-skip}
Let $f$ be an $L$-block residual network with residual blocks $g^{(\ell)}$ and skip connections $h^{(\ell+1)} = h^{(\ell)} + g^{(\ell)}(h^{(\ell)})$. Under Assumptions~\ref{ass:gaussian}, \ref{ass:uniform-width}, \ref{ass:independence}, \ref{ass:smooth}:
\begin{equation}
\label{eq:H-skip}
H_{ijk} = \sum_{\ell=1}^L \Pi_\ell^\top\, H_{ijk}^{(g^{(\ell)})}\, \Pi_\ell + \sum_{\ell=1}^L \Xi_\ell(r, \Theta) + \cO(L^3/N^2),
\end{equation}
where:
\begin{enumerate}[label=(\roman*)]
\item $\Pi_\ell = \prod_{s=\ell+1}^L (I + J_{g^{(s)}})$ is the cumulative Jacobian from layer $\ell$ to the output, with $J_{g^{(s)}} = \pp g^{(s)}/\pp h^{(s)}$.
\item $H_{ijk}^{(g^{(\ell)})}$ is the Hessian--Jacobian tensor of the $\ell$-th residual branch as an independent module.
\item $\Xi_\ell$ collects cross-terms: for $\ell < m$,
\begin{equation}
\Delta H_{ijk}^{(\ell,m)} = \sum_{p,q} \frac{\pp^2 g^{(\ell)}(x_i)}{\pp\theta_p \pp\theta_q}\left[\frac{\pp g^{(\ell)}(x_j)}{\pp\theta_p}\frac{\pp h^{(m)}(x_k)}{\pp\theta_q} + (j \leftrightarrow k)\right],
\end{equation}
mediated by $\Pi_{\ell,m} = \prod_{s=\ell}^{m-1}(I + J_{g^{(s)}})$.
\item At initialization under $\muP$, $\norm{J_{g^{(\ell)}}} = \cO(1/\sqrt{N})$, so $\Pi_\ell = I + \cO(L/\sqrt{N})$ and the cross-terms are $\cO(L/N)$ relative to the direct terms.
\end{enumerate}
\end{theorem}

\begin{proof}
\textbf{Step 1: Jacobian of the composite map.} The full network output is $f = h^{(L)}$ where $h^{(\ell)} = h^{(0)} + \sum_{s=1}^\ell g^{(s)}(h^{(s-1)})$ by telescoping the residual connections. The Jacobian with respect to parameters $\theta^{(\ell)}$ of block $\ell$ is
\begin{equation}
\frac{\pp f}{\pp \theta^{(\ell)}_p} = \Pi_\ell \cdot \frac{\pp g^{(\ell)}}{\pp \theta^{(\ell)}_p},
\end{equation}
where $\Pi_\ell = \prod_{s=\ell+1}^L (I + J_{g^{(s)}})$ propagates the gradient through subsequent skip connections.

\textbf{Step 2: Hessian decomposition.} The second derivative with respect to parameters from blocks $\ell_1$ and $\ell_2$ (with $\ell_1 \le \ell_2$) is:
\begin{align}
\frac{\pp^2 f}{\pp \theta^{(\ell_1)}_p \pp \theta^{(\ell_2)}_q}
&= \begin{cases}
\Pi_{\ell_1} \cdot \frac{\pp^2 g^{(\ell_1)}}{\pp \theta_p \pp \theta_q} & \text{if } \ell_1 = \ell_2, \\[6pt]
\Pi_{\ell_1,\ell_2} \cdot \frac{\pp J_{g^{(\ell_2)}}}{\pp \theta^{(\ell_1)}_p} \cdot \frac{\pp g^{(\ell_2)}}{\pp \theta^{(\ell_2)}_q} \cdot \Pi_{\ell_2} & \text{if } \ell_1 < \ell_2.
\end{cases}
\end{align}
The within-block terms ($\ell_1 = \ell_2$) give $\Pi_\ell^\top H^{(g^{(\ell)})} \Pi_\ell$ upon contraction with Jacobians. The cross-block terms ($\ell_1 < \ell_2$) produce $\Xi_\ell$.

\textbf{Step 3: Norm bounds on cross-terms.} Under $\muP$ with Assumption~\ref{ass:smooth}, the residual branch at initialization satisfies $\norm{g^{(\ell)}(h)} = \cO(1/\sqrt{N})$ and $\norm{J_{g^{(\ell)}}} = \cO(1/\sqrt{N})$. Therefore:
\begin{equation}
\norm{\Pi_\ell - I} \le \sum_{s=\ell+1}^L \norm{J_{g^{(s)}}} + \text{higher-order} = \cO(L/\sqrt{N}).
\end{equation}
The cross-term $\Delta H^{(\ell,m)}$ involves the product $J_{g^{(\ell)}} \cdot J_{g^{(m)}}$ (schematically), which is $\cO(1/N)$. Summing over the $\binom{L}{2}$ pairs gives total cross-term contribution $\cO(L^2/N)$, which is $\cO(L/N)$ relative to the $L$ direct terms.

The remainder from neglected higher-order products of $J_{g^{(\ell)}}$ factors is $\cO(L^3/N^{3/2})$, which for $L \ll N^{1/2}$ is $\cO(L^3/N^2)$ after the overall $1/N$ factor in the expansion.
\end{proof}

\begin{corollary}[Depth scaling]
\label{cor:depth-scaling}
For a depth-$L$ ResNet under $\muP$ with residual branch scaling $g^{(\ell)} \sim 1/\sqrt{L}$:
\begin{equation}
\norm{H_{ijk}} = \cO(1/N),
\end{equation}
uniformly in $L$. Without the $1/\sqrt{L}$ scaling, $\norm{H_{ijk}} = \cO(L/N)$, and the perturbative expansion breaks down at depth $L = \cO(N)$.
\end{corollary}

\begin{proof}
With $g^{(\ell)} \sim 1/\sqrt{L}$, each within-block contribution $H^{(g^{(\ell)})} = \cO(1/(NL))$ (from the $1/\sqrt{L}$ scaling of both the Hessian and Jacobians). Summing $L$ such terms: $\sum_\ell H^{(g^{(\ell)})} = \cO(1/N)$. The cross-terms are further suppressed by $\cO(1/L)$, giving total cross-term $\cO(L^2 \cdot 1/(NL^2)) = \cO(1/N)$.

Without the $1/\sqrt{L}$ scaling, each within-block term is $\cO(1/N)$, and summing $L$ terms gives $\cO(L/N)$.
\end{proof}

\textbf{Implementation mapping.} Theorem~\ref{thm:H-skip} drives the \texttt{ResNetHTensor} module, enabling analytical $H$ computation for residual architectures and informing depth-scaling predictions.

% ---------- 3.4 Kernel Evolution ODE ----------
\subsection{Kernel Evolution ODE}
\label{sec:kernel-ode}

\begin{definition}[Kernel evolution ODE]
\label{def:kernel-ode}
The leading finite-width correction evolves as
\begin{equation}
\label{eq:kernel-ode}
\frac{\dd \Theta^{(1)}_{ij}}{\dd t} = -\eta \sum_{k=1}^n r_k(t)\, H_{ijk}(t),
\end{equation}
coupled with the residual dynamics \eqref{eq:residual-ode}. The full first-order system is:
\begin{equation}
\label{eq:coupled-system}
\begin{cases}
\displaystyle\frac{\dd r_i}{\dd t} = -\eta \sum_j \left[\Theta^{(0)}_{ij} + \frac{1}{N}\Theta^{(1)}_{ij}(t)\right] r_j(t), \\[10pt]
\displaystyle\frac{\dd \Theta^{(1)}_{ij}}{\dd t} = -\eta \sum_k r_k(t)\, H_{ijk}(t),
\end{cases}
\end{equation}
with initial conditions $r_i(0) = f_{\theta_0}(x_i) - y_i$ and $\Theta^{(1)}_{ij}(0) = \Theta_{ij}(\theta_0) - \Theta^{(0)}_{ij}$.
\end{definition}

In principle, $H_{ijk}(t)$ depends on time through the evolving network parameters. In practice, we employ a frozen-$H$ approximation that is justified by the following structural stability result.

\begin{definition}[Frozen-$H$ approximation]
\label{def:frozen-H}
The frozen-$H$ ODE replaces $H_{ijk}(t)$ in \eqref{eq:kernel-ode} with its value at initialization $H_{ijk}(0)$, yielding:
\begin{equation}
\frac{\dd \Theta^{(1)}_{ij}}{\dd t} = -\eta \sum_k r_k(t)\, H_{ijk}(0).
\end{equation}
\end{definition}

\begin{lemma}[Structural stability of frozen-$H$ approximation]
\label{lem:structural-stability}
Let $\gamma_c^{\mathrm{frozen}}$ be the phase boundary predicted by the frozen-$H$ ODE and $\gamma_c^{\mathrm{true}}$ the boundary of the full (time-dependent $H$) system. If $\norm{H(t) - H(0)}_F / \norm{H(0)}_F \le \varepsilon$ for all $t \in [0, T]$, then
\begin{equation}
\label{eq:structural-stability}
|\gamma_c^{\mathrm{frozen}} - \gamma_c^{\mathrm{true}}| \le C \cdot \varepsilon \cdot \gamma_c^{\mathrm{true}},
\end{equation}
where $C$ depends on the spectral gap of $\Theta^{(0)}$ and the training horizon $T$.
\end{lemma}

\begin{proof}
The phase boundary is defined (Definition~\ref{def:phase-boundary}) as the zero-crossing of the maximal Lyapunov exponent $\Lambda_T(\gamma)$, which is a functional of the variational system driven by $H$. We bound the sensitivity of $\Lambda_T$ to perturbations in $H$.

\textbf{Step 1.} The variational system (Definition~\ref{def:phase-boundary}) with $H$ replaced by $H + \delta H$ (where $\delta H = H(t) - H(0)$) has state-transition matrix $\Phi_\varepsilon(T,0)$ satisfying
\begin{equation}
\norm{\Phi_\varepsilon(T,0) - \Phi_0(T,0)} \le \norm{\Phi_0(T,0)} \cdot \left(\exp\!\left(\int_0^T \norm{\delta A(s)}\,\dd s\right) - 1\right),
\end{equation}
where $\delta A(s)$ is the perturbation to the coefficient matrix of the variational system, which is proportional to $\delta H$. Since $\norm{\delta H(t)} \le \varepsilon \norm{H(0)}$, we have $\norm{\delta A(s)} \le C_A \cdot \varepsilon \cdot \gamma$ for a constant $C_A$ depending on $\norm{r_0(t)}$.

\textbf{Step 2.} The Lyapunov exponent is
\begin{equation}
\Lambda_T = \frac{1}{T}\ln \norm{\Phi(T,0)}_{\mathrm{op}}.
\end{equation}
By the matrix perturbation bound,
\begin{equation}
|\Lambda_T^{\mathrm{frozen}} - \Lambda_T^{\mathrm{true}}| \le \frac{C_A \cdot \varepsilon \cdot \gamma}{1} = C_A \varepsilon \gamma.
\end{equation}

\textbf{Step 3.} By Theorem~\ref{thm:stability-boundary} part (ii), $\Lambda_T(\gamma)$ is monotonically increasing in $\gamma$, with $\dd\Lambda_T/\dd\gamma \ge c_0 > 0$ near the phase boundary. Therefore, the shift in the zero-crossing satisfies
\begin{equation}
|\gamma_c^{\mathrm{frozen}} - \gamma_c^{\mathrm{true}}| \le \frac{C_A \varepsilon \gamma_c}{c_0} = C \varepsilon \gamma_c^{\mathrm{true}}.
\end{equation}
\end{proof}

\begin{remark}
The frozen-$H$ approximation is most accurate precisely where it is most needed: near the phase boundary, which lies in the regime where kernel drift is incipient but small. In the deep lazy regime ($\gamma \ll \gamma_c$), kernel drift is negligible by definition ($H(t) \approx H(0)$). In the deep rich regime ($\gamma \gg \gamma_c$), the frozen-$H$ ODE correctly predicts instability even if the quantitative trajectory diverges from the true dynamics. The boundary itself, as the onset of instability, sits where the approximation is best controlled.
\end{remark}

\textbf{Implementation mapping.} The frozen-$H$ ODE is solved by Algorithm~\ref{alg:ode}. The structural stability lemma justifies this approximation and is validated by the frozen-$H$ validation experiment (Section~\ref{sec:frozen-H-validation}).

% ---------- 3.5 Dimensionless Coupling ----------
\subsection{Dimensionless Coupling and $\muP$ Scaling}
\label{sec:gamma}

\begin{definition}[Dimensionless coupling parameter]
\label{def:gamma}
For a network with $\muP$ exponents $(a, b)$ and uniform width $N$:
\begin{equation}
\label{eq:gamma}
\gamma = \eta \cdot N^{-(1-a-b)}.
\end{equation}
\end{definition}

Under standard parameterization (SP, $a=1, b=0$): $\gamma = \eta$. The perturbative expansion parameter is $\gamma/N = \eta/N$, requiring $\eta \to 0$ as $N \to \infty$ for perturbative validity. Under $\muP$ ($a=1, b=1$): $\gamma = \eta N$, and $\gamma = \cO(1)$ at fixed $\eta$, giving a well-ordered expansion at any finite learning rate.

\begin{remark}
The dimensionless coupling $\gamma$ unifies the $\muP$ scaling framework with the perturbative expansion. The phase boundary occurs at a critical value $\gamma_c$ that depends on the architecture, depth, and data, but not on the parameterization convention. This explains why $\muP$ enables hyperparameter transfer: holding $\gamma$ fixed while scaling $N$ keeps the system at the same point in the phase diagram.
\end{remark}

For non-uniform widths (Assumption~\ref{ass:uniform-width} relaxed), the effective coupling is $\gamma_{\mathrm{eff}} = \max_\ell \gamma_\ell$ where $\gamma_\ell = \eta_\ell \cdot N_\ell^{-(1-a_\ell - b_\ell)}$.

% ---------- 3.6 Phase Boundary ----------
\subsection{Phase Boundary as Stability Transition}
\label{sec:phase-boundary}

We now present the corrected formulation of the phase boundary, resolving the fundamental issue with prior steady-state bifurcation approaches.

\begin{remark}[Why steady-state bifurcation fails]
A natural but incorrect approach is to linearize the coupled system \eqref{eq:coupled-system} around the converged state $(\Theta^{(1)} = 0, r^* = 0)$ and look for an eigenvalue crossing. However, at $r^* = 0$, the Jacobian of this system is block-triangular: the kernel perturbation $\delta\Theta$ is driven by the residual through $H_{ijk}$, but the residual equation $\dd(\delta r)/\dd t = -\eta\,\Theta^{(0)}\,\delta r$ has no dependence on $\delta\Theta$ (since the coupling term $-\eta\,\delta\Theta \cdot r^*$ vanishes at $r^* = 0$). Consequently, the eigenvalues are just those of $-\eta\,\Theta^{(0)}$ (always negative) plus zeros from the $\delta\Theta$ block. There is no eigenvalue crossing and no bifurcation.
\end{remark}

The correct formulation considers the stability of the \emph{training trajectory}, not the converged state.

\begin{definition}[Phase boundary via trajectory stability]
\label{def:phase-boundary}
Let $r_0(t)$ be the NTK (frozen-kernel) trajectory solving
\begin{equation}
\frac{\dd r_{0,i}}{\dd t} = -\eta \sum_j \Theta^{(0)}_{ij}\, r_{0,j}(t), \qquad r_0(0) = f_{\theta_0}(x) - y.
\end{equation}
Define the \emph{variational system} for perturbations $(\delta\Theta, \delta r)$ around this trajectory:
\begin{equation}
\label{eq:variational}
\begin{cases}
\displaystyle\frac{\dd(\delta r_i)}{\dd t} = -\eta \sum_j \Theta^{(0)}_{ij}\, \delta r_j - \eta \sum_j \delta\Theta_{ij}\, r_{0,j}(t), \\[10pt]
\displaystyle\frac{\dd(\delta\Theta_{ij})}{\dd t} = -\eta \sum_k r_{0,k}(t)\, H_{ijk}.
\end{cases}
\end{equation}
This is a \emph{non-autonomous linear system} (because $r_0(t)$ is time-dependent). Define the \emph{finite-time Lyapunov exponent}:
\begin{equation}
\label{eq:lyapunov}
\Lambda_T(\gamma) = \frac{1}{T} \ln \frac{\sup_{\norm{v_0}=1} \norm{\Phi(T,0)\, v_0}}{\norm{v_0}},
\end{equation}
where $\Phi(T,0)$ is the state-transition matrix of the variational system \eqref{eq:variational}. The \textbf{phase boundary} is
\begin{equation}
\Gamma = \{(\gamma, L, \cD) : \Lambda_T(\gamma, L, \cD) = 0\}.
\end{equation}
\begin{itemize}
\item $\Lambda_T < 0$: perturbations decay $\to$ frozen-kernel trajectory is stable $\to$ \textbf{lazy regime}.
\item $\Lambda_T > 0$: perturbations grow $\to$ frozen-kernel trajectory is unstable $\to$ \textbf{rich regime}.
\end{itemize}
\end{definition}

\begin{remark}[Feedback mechanism]
The variational system \eqref{eq:variational} is \emph{not} block-triangular. The $\delta r$ equation contains the coupling term $-\eta \sum_j \delta\Theta_{ij}\, r_{0,j}(t)$, which is nonzero whenever $r_0(t) \neq 0$---i.e., during training before convergence. A kernel perturbation $\delta\Theta$ alters the residual evolution via $\delta r$, which in turn drives further kernel evolution through $H_{ijk}$. This genuine two-way feedback loop enables the Lyapunov exponent to cross zero as $\gamma$ increases.
\end{remark}

\begin{definition}[Order parameter: Kernel Alignment Drift Rate (KADR)]
\label{def:KADR}
The instantaneous kernel alignment drift rate is
\begin{equation}
\rho(t) = \frac{\norm{\Theta^{(1)}(t)}_F}{\norm{\Theta^{(0)}}_F},
\end{equation}
and the time-averaged order parameter is $\bar{\rho} = \frac{1}{T}\int_0^T \rho(t)\,\dd t$. The phase is classified as: \textbf{lazy} if $\bar{\rho} < \rho_{\mathrm{lazy}}$ ($\sim 0.1$); \textbf{rich} if $\bar{\rho} > \rho_{\mathrm{rich}}$ ($\sim 0.5$); \textbf{crossover} otherwise.
\end{definition}

\begin{definition}[Operational phase boundary]
\label{def:operational-boundary}
The \emph{primary operational definition} of the phase boundary is the \textbf{steepest-gradient locus of KADR}:
\begin{equation}
\Gamma_{\mathrm{op}} = \left\{\mathbf{p} : \norm{\nabla_{\mathbf{p}} \bar{\rho}(\mathbf{p})} \text{ is locally maximal}\right\},
\end{equation}
where $\mathbf{p} = (\gamma, L, \kappa_{\mathrm{data}})$ collects the control parameters. This definition is robust to finite-$N$ smearing (where the sharp bifurcation becomes a crossover), computable from both theory (via kernel ODE trajectory) and experiment (via training runs), and threshold-free.
\end{definition}

\begin{theorem}[Stability boundary characterization]
\label{thm:stability-boundary}
Consider the variational system \eqref{eq:variational} with uniform width $N$ and dimensionless coupling $\gamma$ (Definition~\ref{def:gamma}). Under Assumptions~\ref{ass:gaussian}--\ref{ass:expansion}, \ref{ass:moment-closure}, \ref{ass:smooth}, and the technical condition $\Theta^{(0)} \succ 0$:
\begin{enumerate}[label=(\roman*)]
\item \textbf{Existence of critical coupling.} There exists $\gamma_c > 0$ such that $\Lambda_T(\gamma) < 0$ for $\gamma < \gamma_c$ and $\Lambda_T(\gamma) > 0$ for $\gamma > \gamma_c$.

\item \textbf{Monotonicity.} $\Lambda_T(\gamma)$ is monotonically increasing in $\gamma$ for fixed $(L, \cD, T)$.

\item \textbf{Agreement with KADR.} The stability boundary $\gamma_c$ satisfies $|\gamma_c - \gamma_c^{\mathrm{KADR}}| = \cO(1/N)$, where $\gamma_c^{\mathrm{KADR}}$ is the steepest-gradient locus of $\bar{\rho}(\gamma)$.

\item \textbf{Finite-width crossover.} At finite $N$, the transition has width $\delta\gamma \sim N^{-\nu}$ with $\nu > 0$, and the order parameter scales as $\bar{\rho} \sim |\gamma - \gamma_c|^\beta$ with $\beta = 1$ at leading order.

\item \textbf{Data dependence.} $\gamma_c$ depends on the data through $\{\lambda_k(\Theta^{(0)})\}$, $\{\ip{v_k}{y}\}$, and the spectrum of $H$, where $\lambda_k, v_k$ are eigenvalues/eigenvectors of $\Theta^{(0)}$.
\end{enumerate}
\end{theorem}

\begin{proof}
\textbf{Part (i): Existence of critical coupling.}

We establish that $\Lambda_T(\gamma)$ transitions from negative to positive as $\gamma$ increases.

\emph{At $\gamma = 0$ (or equivalently, $H = 0$):} The variational system decouples. The $\delta r$ equation becomes $\dd(\delta r)/\dd t = -\eta\,\Theta^{(0)}\,\delta r$, which has Lyapunov exponent $-\eta\,\lambda_{\min}(\Theta^{(0)}) < 0$ since $\Theta^{(0)} \succ 0$. The $\delta\Theta$ equation becomes $\dd(\delta\Theta_{ij})/\dd t = -\eta \sum_k r_{0,k}(t) H_{ijk}$, which is driven by $r_0(t) = e^{-\eta\,\Theta^{(0)}t} r_0(0) \to 0$ exponentially. Thus $\delta\Theta(t)$ converges (bounded), and the combined Lyapunov exponent is $\Lambda_T(0) = -\eta\,\lambda_{\min}(\Theta^{(0)}) < 0$.

\emph{At sufficiently large $\gamma$:} The coupling term $-\eta\,\delta\Theta_{ij}\,r_{0,j}(t)$ in the $\delta r$ equation amplifies perturbations. Consider an initial perturbation $\delta\Theta(0) = H_{ijk} \norm{r_0(0)} / \norm{H}$ (a rank-1 direction aligned with the $H$ tensor). The $\delta\Theta$ equation drives $\delta\Theta$ to grow as $\int_0^t r_0(s)\,H\,\dd s$, which for short times is $\cO(\gamma\,t\,\norm{r_0(0)}\,\norm{H})$. This feeds back into $\delta r$ via the coupling term, producing growth proportional to $\gamma^2 t^2$ in the short-time limit. When this growth rate exceeds the dissipation rate $\eta\,\lambda_{\min}(\Theta^{(0)})$, the net Lyapunov exponent becomes positive. A sufficient condition is:
\begin{equation}
\gamma > \gamma^* \coloneqq \frac{\eta\,\lambda_{\min}(\Theta^{(0)})}{\norm{H}_{\mathrm{op}} \cdot \norm{r_0(0)} \cdot T},
\end{equation}
which is finite and positive.

\emph{Continuity:} The Lyapunov exponent $\Lambda_T(\gamma)$ is a continuous function of $\gamma$ (since the state-transition matrix $\Phi(T,0)$ depends continuously on the coefficient matrix, which depends linearly on $\gamma$). By the intermediate value theorem, there exists $\gamma_c \in (0, \gamma^*)$ with $\Lambda_T(\gamma_c) = 0$.

\textbf{Part (ii): Monotonicity.}

Increasing $\gamma$ (at fixed $H$) amplifies the coupling terms in the variational system by a multiplicative factor. Formally, if $A(t; \gamma)$ is the coefficient matrix of the variational system, then $A(t; \gamma_2) - A(t; \gamma_1) = (\gamma_2 - \gamma_1) \cdot B(t)$ where $B(t)$ has a specific sign structure (the feedback terms have positive effect on perturbation growth). By the comparison theorem for Lyapunov exponents of linear non-autonomous systems: if $A_2(t) \ge A_1(t)$ in the sense that $A_2 - A_1$ is positive semi-definite in the appropriate block structure, then $\Lambda_T(A_2) \ge \Lambda_T(A_1)$.

More precisely, define the growth functional $G(\gamma, v_0) = \ln\norm{\Phi_\gamma(T,0)v_0}$ for unit $v_0$. The coupling terms proportional to $\gamma$ enter with a definite sign in the $(\delta r, \delta\Theta)$ interaction, ensuring that $\pp G/\pp\gamma > 0$ for any $v_0$ that achieves the supremum. Since $\Lambda_T = \sup_{v_0} G/T$, monotonicity follows.

\textbf{Part (iii): Agreement with KADR.}

The KADR order parameter $\bar{\rho}(\gamma) = (1/T)\int_0^T \norm{\Theta^{(1)}(t)}_F/\norm{\Theta^{(0)}}_F\,\dd t$ measures the magnitude of the kernel correction along the actual trajectory of the full system \eqref{eq:coupled-system}. The variational system \eqref{eq:variational} is the linearization of \eqref{eq:coupled-system} around the NTK trajectory; therefore, in the perturbative regime ($1/N \ll 1$), $\Theta^{(1)}(t) = (1/N) \cdot \delta\Theta(t) + \cO(1/N^2)$ where $\delta\Theta$ is the variational perturbation.

The steepest-gradient locus of $\bar{\rho}$ is $\gamma_c^{\mathrm{KADR}} = \argmax_\gamma \norm{\nabla_\gamma \bar{\rho}}$. Near the stability boundary, $\bar{\rho}$ transitions from $\cO(1/N)$ (lazy, bounded $\delta\Theta$) to $\cO(1)$ (rich, exponentially growing $\delta\Theta$). The steepest gradient occurs where $\Lambda_T$ crosses zero, plus an $\cO(1/N)$ correction from the nonlinear terms neglected in the linearization. Therefore $|\gamma_c - \gamma_c^{\mathrm{KADR}}| = \cO(1/N)$.

\textbf{Part (iv): Finite-width crossover.}

At finite $N$, the $1/N$ correction smooths the sharp stability boundary into a crossover. The order parameter satisfies $\bar{\rho}(\gamma) \approx \frac{C}{N} \cdot \frac{|\gamma - \gamma_c|}{\delta\gamma}$ in the crossover region, where $\delta\gamma \sim N^{-\nu}$ is the crossover width. The exponent $\nu$ is determined by the rate at which perturbative corrections sharpen the transition.

The leading-order scaling $\beta = 1$ follows from the linearized analysis: near $\gamma_c$, $\Lambda_T(\gamma) \approx c_0(\gamma - \gamma_c)$ (linear crossing), so $\norm{\delta\Theta(T)} \sim e^{c_0(\gamma - \gamma_c)T}$, yielding $\bar{\rho} \sim |\gamma - \gamma_c|$ (after time-averaging and normalization in the crossover region). Width-dependent corrections from higher-order terms give $\beta_N = 1 + \cO(1/N)$.

\textbf{Part (v): Data dependence.}

The NTK trajectory $r_0(t)$ depends on the data through $\Theta^{(0)}$ and $y$:
\begin{equation}
r_0(t) = e^{-\eta\,\Theta^{(0)}t}\, r_0(0) = \sum_k e^{-\eta\lambda_k t}\, \ip{v_k}{r_0(0)}\, v_k,
\end{equation}
where $\lambda_k, v_k$ are eigenvalues/eigenvectors of $\Theta^{(0)}$ and $r_0(0) = f_{\theta_0}(x) - y$. The variational system's coefficient matrix $A(t;\gamma)$ depends on $r_0(t)$ and $H$, both of which are data-dependent. Thus $\Lambda_T(\gamma)$ and its zero-crossing $\gamma_c$ depend on $\{\lambda_k\}$, $\{\ip{v_k}{y}\}$ (through $r_0(0)$), and the spectral structure of $H$ (which depends on kernel derivatives evaluated on the data).
\end{proof}

\textbf{Implementation mapping.} Theorem~\ref{thm:stability-boundary} drives Algorithm~\ref{alg:growth-rate} (perturbation growth rate computation) and Algorithm~\ref{alg:boundary} (phase boundary detection).

% ---------- 3.7 Moment Closure and UQ ----------
\subsection{Moment Closure and Uncertainty Quantification}
\label{sec:moment-closure}

\begin{definition}[Gaussian moment closure]
\label{def:moment-closure}
The perturbative hierarchy at order $1/N$ requires the 4th cumulant $\kappa_4$ of the pre-activation distribution. The Gaussian closure sets
\begin{equation}
\kappa_4[h^{(\ell)}_{i_1}, h^{(\ell)}_{i_2}, h^{(\ell)}_{i_3}, h^{(\ell)}_{i_4}] = 0
\end{equation}
for all layers $\ell$ and data indices, expressing 4th moments via Isserlis' theorem:
\begin{equation}
\E[h_{i_1}h_{i_2}h_{i_3}h_{i_4}] = K_{i_1i_2}K_{i_3i_4} + K_{i_1i_3}K_{i_2i_4} + K_{i_1i_4}K_{i_2i_3}.
\end{equation}
\end{definition}

\begin{definition}[Perturbative validity indicators]
\label{def:validity}
A prediction at parameters $(\gamma, N, L)$ is assessed via:
\begin{enumerate}[label=(\arabic*)]
\item \textbf{Expansion ratio:} $\epsilon_1 = \sup_{t \in [0,T]} \norm{\Theta^{(1)}(t)}_F / (N \cdot \norm{\Theta^{(0)}}_F)$. Require $\epsilon_1 < 0.5$.
\item \textbf{Moment-closure quality:} $\epsilon_2 = \norm{\kappa_4}_{\mathrm{eff}} / \norm{\kappa_2}_{\mathrm{eff}}^2$, estimated from Monte Carlo samples of pre-activations at initialization. Require $\epsilon_2 < 0.3$.
\item \textbf{Two-loop convergence:} $\epsilon_3 = |\gamma_c^{(1\text{-loop})} - \gamma_c^{(2\text{-loop})}| / \gamma_c^{(1\text{-loop})}$. Require $\epsilon_3 < 0.15$.
\item \textbf{Nystr\"om error:} $\epsilon_4 = \norm{\Theta^{(0)}_m - \Theta^{(0)}}_F / \norm{\Theta^{(0)}}_F$ for rank-$m$ approximation. Require $\epsilon_4 < 0.1$.
\end{enumerate}
Confidence classification: \textbf{High} if all $\epsilon_i$ below thresholds; \textbf{Moderate} if at most one marginal; \textbf{Low} if any exceeds threshold.
\end{definition}

\begin{theorem}[Perturbative validity bounds]
\label{thm:validity}
Under Assumptions~\ref{ass:gaussian}--\ref{ass:uniform-width}, \ref{ass:independence}--\ref{ass:horizon}, the one-loop correction satisfies:
\begin{equation}
\label{eq:validity-bound}
\sup_{t \in [0,T]} \frac{\norm{\Theta^{(1)}(t)}_F}{N \cdot \norm{\Theta^{(0)}}_F} \le \frac{C_1 \cdot \gamma \cdot \norm{y}^2}{N \cdot \lambda_{\min}(\Theta^{(0)})} \cdot \left(e^{C_2 \eta \norm{\Theta^{(0)}}_{\mathrm{op}} T} - 1\right),
\end{equation}
where $C_1$ depends on $L$, $\norm{\sigma'}_\infty$, $\norm{\sigma''}_\infty$, and $C_2 = \cO(L)$.
\end{theorem}

\begin{proof}
\textbf{Step 1: Bound on $H_{ijk}$.} From the layer-wise decomposition (Definition~\ref{def:H-mlp}), each layer contributes $H^{(\ell)}_{ijk}$ with
\begin{equation}
|H^{(\ell)}_{ijk}| \le \frac{C_H(\sigma)}{N_\ell} \cdot \sup_{x,x'} |K^{(\ell)}(x,x')| \cdot \sup_{x,x'} |\dot{K}^{(\ell)}(x,x')| \cdot \norm{\Theta^{(\ell+1:L)}}_{\mathrm{op}}^2,
\end{equation}
where $C_H(\sigma)$ depends on $\norm{\sigma'}_\infty$ and $\norm{\sigma''}_\infty$ (bounded by Assumption~\ref{ass:smooth}). The kernel norms are bounded at initialization by constants depending on $L$ and $\sigma$. Summing over $L$ layers:
\begin{equation}
\norm{H}_{\mathrm{op}} \le \frac{L \cdot C_H}{N} \eqqcolon \frac{C_H^{\mathrm{tot}}}{N}.
\end{equation}

\textbf{Step 2: Residual bound.} From the coupled system \eqref{eq:coupled-system}, using the perturbative ordering $\Theta^{(1)}/N \ll \Theta^{(0)}$:
\begin{equation}
\norm{r(t)} \le \norm{r(0)} \exp\!\left(-\eta\,\lambda_{\min}(\Theta^{(0)})\,t + \frac{\eta}{N}\int_0^t \norm{\Theta^{(1)}(s)}_{\mathrm{op}}\,\dd s\right).
\end{equation}

\textbf{Step 3: Growth bound on $\Theta^{(1)}$.} From \eqref{eq:kernel-ode}:
\begin{equation}
\frac{\dd}{\dd t}\norm{\Theta^{(1)}}_F \le \eta\,\norm{r(t)} \cdot \norm{H}_{\mathrm{op}} \cdot n \le \eta\,n\,\norm{r(0)}\,\frac{C_H^{\mathrm{tot}}}{N}\,e^{-\eta\lambda_{\min}(\Theta^{(0)})t + \eta\int_0^t \norm{\Theta^{(1)}(s)}_{\mathrm{op}}/N\,\dd s}.
\end{equation}

\textbf{Step 4: Gr\"onwall inequality.} Define $u(t) = \norm{\Theta^{(1)}(t)}_F$. The differential inequality has the form $\dot{u} \le a(t) + b(t)\,u$ where $b(t)$ captures the feedback of $\Theta^{(1)}$ on $r$. Applying Gr\"onwall's lemma:
\begin{equation}
u(t) \le u(0)\,e^{\int_0^t b(s)\,\dd s} + \int_0^t a(s)\,e^{\int_s^t b(\tau)\,\dd\tau}\,\dd s.
\end{equation}
The dominant term in the exponent is $\eta\,\norm{\Theta^{(0)}}_{\mathrm{op}} \cdot t$ (from the residual's interaction with the kernel), yielding the exponential factor in \eqref{eq:validity-bound}. The prefactor $C_1 \gamma \norm{y}^2 / (N\lambda_{\min})$ collects the initial residual norm ($\norm{r(0)} \le \norm{y}$ at zero-initialized output), the coupling strength, and the spectral conditioning.
\end{proof}

\begin{remark}[Conservativeness of the bound]
The exponential factor $e^{C_2 \eta \norm{\Theta^{(0)}} T}$ makes this bound vacuous for practical training durations. For example, with $\eta = 0.1$, $\norm{\Theta^{(0)}} \sim n = 200$, $T = 2000$, and $C_2 \sim L = 4$, the exponent is $\sim 10^5$. We emphasize that Theorem~\ref{thm:validity} provides \emph{qualitative} guidance (the expansion degrades with increasing $\gamma$, $T$, and conditioning) while the \emph{quantitative} diagnostic is the computable indicator $\epsilon_1$ evaluated on actual ODE trajectories. The bound establishes that perturbative validity is a well-defined concept with a formal sufficient condition.
\end{remark}

\textbf{Implementation mapping.} Theorem~\ref{thm:validity} motivates the \texttt{ValidityChecker} module (Algorithm~\ref{alg:confidence}), which computes $\epsilon_1$ from ODE output rather than relying on the conservative Gr\"onwall bound.

% ---------- 3.8 ReLU Extension ----------
\subsection{ReLU Extension}
\label{sec:relu}

All theorems in this paper require Assumption~\ref{ass:smooth} ($C^3$ activation). ReLU does not satisfy this assumption: its derivative is discontinuous at zero, and its second derivative is a Dirac delta in the distributional sense.

\begin{conjecture}[ReLU extension]
\label{conj:relu}
Theorems~\ref{thm:expansion}--\ref{thm:validity} hold for ReLU activation with the following modifications:
\begin{enumerate}[label=(\roman*)]
\item $\Theta^{(0)}$ is replaced by the arc-cosine kernel \citep{cho2009kernel}.
\item The $C^3$ bounds on $\sigma$ are replaced by distributional estimates using the piecewise-linear structure of ReLU.
\item Phase boundary locations and scaling exponents are qualitatively preserved.
\end{enumerate}
\end{conjecture}

\begin{remark}
We test Conjecture~\ref{conj:relu} empirically by running the identical experimental protocol with both GELU (primary, satisfying Assumption~\ref{ass:smooth}) and ReLU. The comparison provides evidence for or against the conjecture. For ReLU, we additionally validate that the extrapolated $\Theta^{(0)}$ from calibration widths converges to the known arc-cosine kernel, providing a consistency check on the expansion.
\end{remark}

% =============================================================================
\section{Algorithms}
\label{sec:algorithms}

We present the complete algorithmic pipeline, from NTK computation through phase diagram construction. Each algorithm includes pseudocode, a complexity analysis, and references to the theorems it implements.

\subsection{Exact NTK Computation}
\label{sec:alg-ntk}

\begin{algorithm}[H]
\caption{Exact NTK Computation via Autodiff}
\label{alg:ntk}
\DontPrintSemicolon
\KwIn{Network $f(\cdot;\theta)$, dataset $X = \{x_1,\ldots,x_n\}$, parameters $\theta$}
\KwOut{NTK matrix $\Theta \in \R^{n \times n}$, per-layer decomposition $\{\Theta^{(\ell)}\}_{\ell=1}^L$}
\BlankLine
$J \gets \mathrm{zeros}(n, P)$ \tcp*{Jacobian matrix}
\For{$i = 1$ \KwTo $n$}{
  $J[i,:] \gets \nabla_\theta f(x_i; \theta)$ \tcp*{reverse-mode autodiff}
}
\For{$\ell = 1$ \KwTo $L$}{
  $J_\ell \gets J[:, \mathrm{param\_range}(\ell)]$ \tcp*{slice by layer}
  $\Theta^{(\ell)} \gets J_\ell J_\ell^\top$ \tcp*{per-layer NTK}
}
$\Theta \gets \sum_\ell \Theta^{(\ell)}$\;
\Return $\Theta, \{\Theta^{(\ell)}\}_{\ell=1}^L$
\end{algorithm}

\textbf{Complexity.} Time: $\cO(n \cdot P + n^2 \cdot P)$ where $P = \cO(LN^2)$ is the parameter count. For the Jacobian computation, each backward pass costs $\cO(P)$; the outer product is $\cO(n^2 P)$. Space: $\cO(nP + n^2)$ for the Jacobian and kernel matrix.

\textbf{Correctness.} By Definition~\ref{def:ntk}, $\Theta_{ij} = \sum_p (\pp f_i/\pp\theta_p)(\pp f_j/\pp\theta_p) = (JJ^\top)_{ij}$. Vectorized implementations use \texttt{vmap} over data points for efficiency.

\subsection{Finite-Width Correction Extraction}
\label{sec:alg-correction}

\begin{algorithm}[H]
\caption{Finite-Width Correction Extraction}
\label{alg:correction}
\DontPrintSemicolon
\KwIn{Widths $\{N_1,\ldots,N_W\}$ ($W \ge 3$), $S$ kernel samples per width, dataset $X$}
\KwOut{$\Theta^{(0)}, \Theta^{(1)}$, standard errors $\sigma_0, \sigma_1$, residual $R$}
\BlankLine
\tcp{Average over seeds at each width}
\For{$w = 1$ \KwTo $W$}{
  $\bar{\Theta}_w \gets \frac{1}{S}\sum_{s=1}^S \Theta^{(s)}(N_w)$ \;
  $V_w \gets \frac{1}{S-1}\sum_{s=1}^S (\Theta^{(s)}(N_w) - \bar{\Theta}_w)^{\odot 2}$ \tcp*{element-wise variance}
}
\tcp{Centered WLS regression per entry}
$\bar{x} \gets \frac{1}{W}\sum_w 1/N_w$ \tcp*{centering for conditioning}
\For{$i = 1$ \KwTo $n$, $j = i$ \KwTo $n$}{
  $y_{\mathrm{vec}} \gets [\bar{\Theta}_1[i,j], \ldots, \bar{\Theta}_W[i,j]]^\top$ \;
  $x_{\mathrm{vec}} \gets [1/N_1 - \bar{x}, \ldots, 1/N_W - \bar{x}]^\top$ \tcp*{centered regressor}
  $w_{\mathrm{vec}} \gets [S/V_1[i,j], \ldots, S/V_W[i,j]]^\top$ \tcp*{inverse-variance weights}
  \BlankLine
  $A \gets [\mathbf{1}, x_{\mathrm{vec}}]$ \tcp*{$W \times 2$ design matrix}
  $\hat{\beta} \gets (A^\top \mathrm{diag}(w_{\mathrm{vec}}) A)^{-1} A^\top \mathrm{diag}(w_{\mathrm{vec}}) y_{\mathrm{vec}}$ \;
  $\Theta^{(0)}[i,j] \gets \hat{\beta}_0 - \hat{\beta}_1 \bar{x}$ \tcp*{de-center intercept}
  $\Theta^{(1)}[i,j] \gets \hat{\beta}_1$ \;
  $\mathrm{Cov}_\beta \gets (A^\top \mathrm{diag}(w_{\mathrm{vec}}) A)^{-1}$ \;
  $\sigma_0[i,j] \gets \sqrt{\mathrm{Cov}_\beta[0,0]}$;\quad $\sigma_1[i,j] \gets \sqrt{\mathrm{Cov}_\beta[1,1]}$ \;
}
Symmetrize $\Theta^{(0)}, \Theta^{(1)}, \sigma_0, \sigma_1$ \;
$R \gets \sqrt{\sum_{i \le j} \norm{y_{\mathrm{vec}} - A\hat{\beta}}^2} / \binom{n+1}{2}$ \;
\If{$\max|\Theta^{(1)}| / \max|\Theta^{(0)}| > 0.5$}{
  \textbf{warn} ``Perturbative expansion may not converge'' \;
}
\Return $\Theta^{(0)}, \Theta^{(1)}, \sigma_0, \sigma_1, R$
\end{algorithm}

\textbf{Complexity.} Time: $\cO(W \cdot S \cdot T_{\mathrm{NTK}} + n^2 W)$. The NTK computations dominate. The regression per entry is $\cO(W)$. Space: $\cO(W S n^2)$.

\textbf{Correctness.} Implements the model $\E[\Theta(N)] = \Theta^{(0)} + (1/N)\Theta^{(1)} + \cO(1/N^2)$ from Theorem~\ref{thm:expansion}. Centering the regressor (replacing $1/N$ with $1/N - \overline{1/N}$) reduces the condition number of the normal equation matrix from $\cO(N_{\max}^2) \sim 10^6$ to $\cO(N_{\max})$, addressing the ill-conditioning identified in Section~\ref{sec:limitations}. Bootstrap confidence intervals over seeds complement the analytical standard errors.

\subsection{Kernel Evolution ODE Solver}
\label{sec:alg-ode}

\begin{algorithm}[H]
\caption{Kernel Evolution ODE with Frozen-$H$}
\label{alg:ode}
\DontPrintSemicolon
\KwIn{$\Theta^{(0)}, \Theta^{(1)}(0), H, y, \eta, N, T, \epsilon_{\mathrm{tol}}$}
\KwOut{Trajectories $\{(\Theta(t_k), r(t_k))\}_{k=0}^K$, convergence flag}
\BlankLine
$\Theta_{\mathrm{init}} \gets \Theta^{(0)} + \Theta^{(1)}(0)/N$ \;
$r_0 \gets -y$ \tcp*{$f_{\theta_0} = 0$ WLOG}
$z_0 \gets [\mathrm{upper\_tri}(\Theta_{\mathrm{init}}); r_0]$ \;
\BlankLine
\SetKwFunction{RHS}{RHS}
\SetKwProg{Fn}{Function}{:}{}
\Fn{\RHS{$t, z$}}{
  $\Theta \gets \mathrm{unflatten}(z[1:D_\Theta])$;\quad $r \gets z[D_\Theta+1:\mathrm{end}]$ \;
  $\dd r/\dd t \gets -\eta\,\Theta\, r$ \;
  $\dd\Theta/\dd t \gets -(\eta/N)\,\mathrm{einsum}(\text{`ijk,k$\to$ij'}, H, r)$ \tcp*{frozen $H$}
  \Return $[\mathrm{upper\_tri}(\dd\Theta/\dd t); \dd r/\dd t]$
}
\BlankLine
\tcp{Stiffness detection}
$\lambda_{\max} \gets \norm{\Theta^{(0)}}_{\mathrm{op}}$ \;
\eIf{$\eta \cdot \lambda_{\max} \cdot T > 1000$}{
  method $\gets$ BDF \tcp*{implicit, for stiff systems}
}{
  method $\gets$ RK45 \tcp*{explicit adaptive}
}
$(t_{\mathrm{pts}}, z_{\mathrm{traj}}) \gets \mathrm{ODE\_solve}(\mathrm{RHS}, z_0, [0,T], \text{method}, \epsilon_{\mathrm{tol}})$ \;
converged $\gets \norm{r(T)} < \epsilon_{\mathrm{tol}} \cdot \norm{r(0)}$ \;
\Return $(t_{\mathrm{pts}}, \Theta_{\mathrm{traj}}, r_{\mathrm{traj}}, \text{converged})$
\end{algorithm}

\textbf{Complexity.} Time: $\cO(K_{\mathrm{steps}} \cdot n^3)$ per grid point, where the $H$-tensor contraction (einsum) dominates each RHS evaluation. Space: $\cO(n^3)$ for $H$ storage plus $\cO(n^2)$ for the state.

\textbf{Correctness.} Implements Definition~\ref{def:frozen-H}. The frozen-$H$ approximation is justified by Lemma~\ref{lem:structural-stability}. Adaptive step-size control (RK45 or BDF) ensures the ODE discretization error is below $\epsilon_{\mathrm{tol}}$.

\subsection{Perturbation Growth Rate Computation}
\label{sec:alg-growth}

\begin{algorithm}[H]
\caption{Finite-Time Lyapunov Exponent via QR Decomposition}
\label{alg:growth-rate}
\DontPrintSemicolon
\KwIn{$\Theta^{(0)}, H, r_0(t)$ (NTK trajectory), $\gamma, \eta, T, k$ (number of exponents)}
\KwOut{Lyapunov exponents $\{\Lambda_1, \ldots, \Lambda_k\}$}
\BlankLine
$D \gets n(n+1)/2 + n$ \tcp*{state dimension}
$Q \gets \mathrm{orthogonal\_basis}(D, k)$ \tcp*{$D \times k$ initial basis}
$R_{\mathrm{accum}} \gets I_k$ \tcp*{accumulated $R$ factors}
$\Delta t \gets T / M$ \tcp*{$M$ reorthogonalization intervals}
\BlankLine
\SetKwFunction{Amat}{VariationalRHS}
\Fn{\Amat{$t, v$}}{
  $v_\Theta \gets v[1:D_\Theta]$;\quad $v_r \gets v[D_\Theta+1:\mathrm{end}]$ \;
  $w_r \gets -\eta\,\Theta^{(0)}\,v_r - \eta\gamma\,\mathrm{einsum}(\text{`ij,j$\to$i'},\mathrm{unflatten}(v_\Theta), r_0(t))$ \;
  $w_\Theta \gets -\eta\gamma\,\mathrm{einsum}(\text{`ijk,k$\to$ij'}, H, v_r) \cdot \norm{r_0(t)}$ \;
  \Return $[w_\Theta; w_r]$
}
\BlankLine
\For{$m = 1$ \KwTo $M$}{
  \tcp{Evolve $k$ columns of $Q$ over interval $[(m-1)\Delta t, m\Delta t]$}
  \For{$j = 1$ \KwTo $k$}{
    $Q[:,j] \gets \mathrm{ODE\_step}(\mathrm{VariationalRHS}, Q[:,j], [(m\!-\!1)\Delta t, m\Delta t])$
  }
  \tcp{QR reorthogonalization to prevent collapse}
  $(Q, R) \gets \mathrm{QR}(Q)$ \tcp*{thin QR: $D \times k$}
  $R_{\mathrm{accum}} \gets R_{\mathrm{accum}} \cdot R$ \;
}
\BlankLine
\For{$j = 1$ \KwTo $k$}{
  $\Lambda_j \gets \frac{1}{T} \sum_{m=1}^M \ln |R_m[j,j]|$ \tcp*{sum of log diagonal entries}
}
\Return $\{\Lambda_1, \ldots, \Lambda_k\}$ \tcp*{sorted by magnitude}
\end{algorithm}

\textbf{Complexity.} Time: $\cO(M \cdot k \cdot T_{\mathrm{step}} + M \cdot D \cdot k^2)$ where $T_{\mathrm{step}}$ is the cost of one ODE step of the variational system. The variational RHS is $\cO(n^3)$ (dominated by the $H$-contraction). With $M \sim 100$ reorthogonalization intervals and $k \sim 10$ exponents: $\cO(100 \cdot 10 \cdot n^3 + 100 \cdot n^2 \cdot 100) \approx \cO(1000\,n^3)$. Space: $\cO(Dk + n^3)$.

\textbf{Correctness.} The QR-based algorithm for Lyapunov exponents is standard in dynamical systems \citep{benettin1980lyapunov}. Periodic QR reorthogonalization prevents the $k$ basis vectors from collapsing onto the most unstable direction, enabling computation of sub-dominant exponents. The accumulated diagonal entries of $R$ track the growth rates, converging to the true Lyapunov exponents as $T \to \infty$.

\textbf{Why QR, not Arnoldi.} The variational system is non-autonomous ($r_0(t)$ is time-dependent), so the coefficient matrix $A(t)$ changes at each time step. Standard Arnoldi methods target eigenvalues of a fixed matrix and are not directly applicable. The QR-based Lyapunov exponent algorithm handles non-autonomous systems natively.

\subsection{Phase Boundary Detection}
\label{sec:alg-boundary}

\begin{algorithm}[H]
\caption{Phase Boundary Detection via Bisection on Growth Rate}
\label{alg:boundary}
\DontPrintSemicolon
\KwIn{Growth rate oracle $\Lambda_T(\gamma)$, bracket $[\gamma_{\mathrm{lo}}, \gamma_{\mathrm{hi}}]$, tolerance $\epsilon_{\mathrm{bif}}$}
\KwOut{Critical coupling $\gamma_c$}
\BlankLine
\tcp{Verify bracket}
$\Lambda_{\mathrm{lo}} \gets \Lambda_T(\gamma_{\mathrm{lo}})$;\quad $\Lambda_{\mathrm{hi}} \gets \Lambda_T(\gamma_{\mathrm{hi}})$ \;
\textbf{assert} $\Lambda_{\mathrm{lo}} < 0$ and $\Lambda_{\mathrm{hi}} > 0$ \;
\BlankLine
\tcp{Bisection phase}
\While{$\gamma_{\mathrm{hi}} - \gamma_{\mathrm{lo}} > \epsilon_{\mathrm{bif}}$}{
  $\gamma_{\mathrm{mid}} \gets (\gamma_{\mathrm{lo}} + \gamma_{\mathrm{hi}})/2$ \;
  $\Lambda_{\mathrm{mid}} \gets \Lambda_T(\gamma_{\mathrm{mid}})$ \;
  \eIf{$\Lambda_{\mathrm{mid}} < 0$}{
    $\gamma_{\mathrm{lo}} \gets \gamma_{\mathrm{mid}}$;\quad $\Lambda_{\mathrm{lo}} \gets \Lambda_{\mathrm{mid}}$
  }{
    $\gamma_{\mathrm{hi}} \gets \gamma_{\mathrm{mid}}$;\quad $\Lambda_{\mathrm{hi}} \gets \Lambda_{\mathrm{mid}}$
  }
}
\tcp{Secant refinement}
\For{$\mathrm{iter} = 1$ \KwTo $5$}{
  $\gamma_{\mathrm{new}} \gets \gamma_{\mathrm{lo}} - \Lambda_{\mathrm{lo}} \cdot \frac{\gamma_{\mathrm{hi}} - \gamma_{\mathrm{lo}}}{\Lambda_{\mathrm{hi}} - \Lambda_{\mathrm{lo}}}$ \;
  $\gamma_{\mathrm{new}} \gets \mathrm{clamp}(\gamma_{\mathrm{new}}, \gamma_{\mathrm{lo}}, \gamma_{\mathrm{hi}})$ \;
  $\Lambda_{\mathrm{new}} \gets \Lambda_T(\gamma_{\mathrm{new}})$ \;
  \lIf{$|\Lambda_{\mathrm{new}}| < \epsilon_{\mathrm{bif}} \cdot |\Lambda_{\mathrm{lo}}|$}{\textbf{break}}
  \eIf{$\Lambda_{\mathrm{new}} < 0$}{$(\gamma_{\mathrm{lo}}, \Lambda_{\mathrm{lo}}) \gets (\gamma_{\mathrm{new}}, \Lambda_{\mathrm{new}})$}{$(\gamma_{\mathrm{hi}}, \Lambda_{\mathrm{hi}}) \gets (\gamma_{\mathrm{new}}, \Lambda_{\mathrm{new}})$}
}
$\gamma_c \gets (\gamma_{\mathrm{lo}} + \gamma_{\mathrm{hi}})/2$ \;
\Return $\gamma_c$
\end{algorithm}

\textbf{Complexity.} Time: $\cO(\log((\gamma_{\mathrm{hi}} - \gamma_{\mathrm{lo}})/\epsilon_{\mathrm{bif}}) \cdot T_{\Lambda})$ where $T_\Lambda$ is the cost of one Lyapunov exponent computation ($\cO(1000\,n^3)$ from Algorithm~\ref{alg:growth-rate}). With $\epsilon_{\mathrm{bif}} = 10^{-6}$ and bracket width $\sim 1$: $\sim 20$ bisection steps plus $\le 5$ secant refinements.

\textbf{Correctness.} By Theorem~\ref{thm:stability-boundary}(i)--(ii), $\Lambda_T(\gamma)$ is continuous and monotonically increasing, so bisection is guaranteed to converge. The secant method provides superlinear refinement.

\subsection{Pseudo-Arclength Continuation}
\label{sec:alg-continuation}

\begin{algorithm}[H]
\caption{Pseudo-Arclength Continuation for Phase Boundaries}
\label{alg:continuation}
\DontPrintSemicolon
\KwIn{$G(\gamma, p) = \Lambda_T(\gamma, p)$, initial point $(\gamma_0, p_0)$ with $G = 0$, step $\Delta s$, bound $s_{\max}$}
\KwOut{Boundary curve $\{(\gamma_k, p_k)\}_{k=0}^K$}
\BlankLine
curve $\gets [(\gamma_0, p_0)]$ \;
$\tau \gets \mathrm{normalize}(-\pp G/\pp p, \pp G/\pp\gamma)$ \tcp*{tangent via finite diff.}
$s \gets 0$ \;
\While{$s < s_{\max}$}{
  \tcp{Predictor: Euler step along tangent}
  $(\gamma_{\mathrm{pred}}, p_{\mathrm{pred}}) \gets (\gamma_k, p_k) + \Delta s \cdot \tau$ \;
  \tcp{Corrector: Newton on augmented system}
  $(\gamma_{k+1}, p_{k+1}) \gets (\gamma_{\mathrm{pred}}, p_{\mathrm{pred}})$ \;
  \For{$\mathrm{iter} = 1$ \KwTo $20$}{
    $g \gets G(\gamma_{k+1}, p_{k+1})$ \;
    $\nabla g \gets (\pp G/\pp\gamma, \pp G/\pp p)$ \tcp*{finite differences}
    $h \gets \tau \cdot [(\gamma_{k+1}, p_{k+1}) - (\gamma_{\mathrm{pred}}, p_{\mathrm{pred}})]$ \;
    $J_{\mathrm{aug}} \gets \begin{bmatrix} \nabla g \\ \tau^\top \end{bmatrix}$ \;
    $(\Delta\gamma, \Delta p) \gets J_{\mathrm{aug}}^{-1} (-g, -h)$ \;
    $(\gamma_{k+1}, p_{k+1}) \mathrel{+}= (\Delta\gamma, \Delta p)$ \;
    \lIf{$\sqrt{\Delta\gamma^2 + \Delta p^2} < 10^{-10}$}{\textbf{break}}
  }
  \lIf{$|G(\gamma_{k+1}, p_{k+1})| > 10^{-6}$}{$\Delta s \gets \Delta s/2$; \textbf{continue}}
  \tcp{Update tangent and adapt step size}
  $\tau_{\mathrm{new}} \gets \mathrm{normalize}(-\pp G/\pp p, \pp G/\pp\gamma)$ \;
  \lIf{$\tau_{\mathrm{new}} \cdot \tau < 0$}{$\tau_{\mathrm{new}} \gets -\tau_{\mathrm{new}}$}
  $\cos\alpha \gets \tau_{\mathrm{new}} \cdot \tau$ \;
  \lIf{$\cos\alpha < 0.5$}{$\Delta s \gets \max(\Delta s/2, \Delta s_{\min})$}
  \lElseIf{$\cos\alpha > 0.95$}{$\Delta s \gets \min(1.5\Delta s, \Delta s_{\max})$}
  curve.append$(\gamma_{k+1}, p_{k+1})$;\quad $\tau \gets \tau_{\mathrm{new}}$;\quad $s \mathrel{+}= \Delta s$ \;
}
\Return curve
\end{algorithm}

\textbf{Complexity.} Each continuation step requires $\cO(n_{\mathrm{Newton}} \cdot T_G)$ where $T_G = T_\Lambda$ (one growth rate evaluation per $G$ call) and the gradient requires 2 additional evaluations (central finite differences). Total: $\cO(K_{\mathrm{steps}} \cdot 60 \cdot T_\Lambda)$ accounting for Newton iterations and gradient evaluations.

\textbf{Correctness.} Pseudo-arclength continuation \citep{keller1977numerical} avoids the singularity at fold points (where $\pp G/\pp\gamma = 0$) by augmenting the constraint equation with an arclength condition. The tangent update ensures consistent orientation. Adaptive step-size control based on the angle between successive tangents prevents the corrector from diverging at high-curvature points.

\subsection{Nystr\"om Approximation}
\label{sec:alg-nystrom}

\begin{algorithm}[H]
\caption{Nystr\"om Kernel Approximation with Adaptive Rank}
\label{alg:nystrom}
\DontPrintSemicolon
\KwIn{Dataset $X$, kernel function $k(\cdot,\cdot)$, target rank $m$, strategy $\in \{\text{uniform, k-means, leverage}\}$}
\KwOut{Low-rank factor $L \in \R^{n \times r}$ such that $\Theta \approx LL^\top$, error $\epsilon_{\mathrm{nys}}$}
\BlankLine
Select $m$ landmark indices by chosen strategy \;
$K_{mm} \gets k(X_m, X_m) \in \R^{m \times m}$ \;
$K_{nm} \gets k(X, X_m) \in \R^{n \times m}$ \;
$(U, \Sigma, V^\top) \gets \mathrm{SVD}(K_{mm})$ \;
$r \gets |\{i : \Sigma_i > 10^{-10} \cdot \Sigma_1\}|$ \tcp*{effective rank}
$L \gets K_{nm} \cdot U_{:,:r} \cdot \mathrm{diag}(1/\sqrt{\Sigma_{1:r}})$ \tcp*{$n \times r$ factor}
\tcp{Error estimation via subsampling}
$\epsilon_{\mathrm{nys}} \gets \norm{k(X_{\mathrm{check}}, X_{\mathrm{check}}) - L_{\mathrm{check}} L_{\mathrm{check}}^\top}_F / \norm{k(X_{\mathrm{check}}, X_{\mathrm{check}})}_F$ \;
\Return $L, \epsilon_{\mathrm{nys}}, r$
\end{algorithm}

\textbf{Complexity.} Time: $\cO(nm \cdot T_k + m^3)$ where $T_k$ is the cost of one kernel evaluation. Space: $\cO(nr + m^2)$, replacing $\cO(n^2)$ for the full kernel.

\textbf{Implementation note.} At our primary probe set size $n_{\mathrm{probe}} = 200$, Nystr\"om is unnecessary (exact kernel computation is tractable). Nystr\"om is reserved for scaling experiments with $n > 1000$.

\subsection{Phase Diagram Orchestration}
\label{sec:alg-orchestration}

\begin{algorithm}[H]
\caption{Phase Diagram Sweep Pipeline}
\label{alg:sweep}
\DontPrintSemicolon
\KwIn{Architecture spec, dataset $(X,y)$, parameter grid $\eta \times N \times \ldots$}
\KwOut{Phase diagram (regime labels, boundaries, confidence, order parameters)}
\BlankLine
\tcp{Stage 1: Amortized kernel computation}
\For{$N_{\mathrm{cal}} \in \{64, 128, 256, 512, 1024\}$}{
  \For{$s = 1$ \KwTo $S$}{
    $\Theta^{(s)}(N_{\mathrm{cal}}) \gets \textsc{ComputeNTK}(f, X, \theta^{(s)}_{N_{\mathrm{cal}}})$ \tcp*{Alg.~\ref{alg:ntk}}
  }
}
$(\Theta^{(0)}, \Theta^{(1)}, \sigma_0, \sigma_1, R) \gets \textsc{ExtractCorrection}(\ldots)$ \tcp*{Alg.~\ref{alg:correction}}
$H \gets \textsc{ComputeHTensor}(\text{arch}, X, \text{cal.\ widths}, \text{kernels})$ \;
\BlankLine
\tcp{Stage 2: Per-grid-point classification}
\ForPar{$(N, \eta) \in \text{grid}$}{
  $\gamma \gets \eta \cdot N^{-(1-a-b)}$ \;
  $(\Theta_{\mathrm{traj}}, r_{\mathrm{traj}}) \gets \textsc{SolveODE}(\Theta^{(0)}, \Theta^{(1)}, H, y, \eta, N, T)$ \tcp*{Alg.~\ref{alg:ode}}
  $\Lambda_1 \gets \textsc{LyapunovExponent}(\Theta^{(0)}, H, r_0, \gamma, \eta, T)$ \tcp*{Alg.~\ref{alg:growth-rate}}
  $\bar{\rho} \gets \textsc{ComputeKADR}(\Theta_{\mathrm{traj}}, \Theta^{(0)})$ \;
  regime $\gets$ classify($\Lambda_1, \bar{\rho}$) \;
  confidence $\gets \textsc{Confidence}(\epsilon_1, \epsilon_2, \epsilon_3, \epsilon_4)$ \tcp*{Alg.~\ref{alg:confidence}}
}
\BlankLine
\tcp{Stage 3: Boundary detection and continuation}
\For{adjacent pairs $(p_1, p_2)$ with different regimes}{
  $\gamma_c \gets \textsc{DetectBoundary}(\Lambda_T, \gamma_{\mathrm{lo}}, \gamma_{\mathrm{hi}})$ \tcp*{Alg.~\ref{alg:boundary}}
  boundary $\gets \textsc{Continue}(\Lambda_T, \gamma_c, p, \Delta s)$ \tcp*{Alg.~\ref{alg:continuation}}
}
\Return phase diagram
\end{algorithm}

\subsection{Order Parameter Computation}
\label{sec:alg-kadr}

The KADR (Definition~\ref{def:KADR}) is computed from the ODE trajectory as:
\begin{equation}
\bar{\rho} = \frac{1}{K}\sum_{k=1}^K \frac{\norm{\Theta_{\mathrm{traj}}(t_k) - \Theta^{(0)}}_F}{\norm{\Theta^{(0)}}_F}.
\end{equation}
For ground-truth measurements, KADR is computed from actual training runs using the empirical NTK recomputed every $\Delta t = 50$ gradient steps on a fixed probe set of 200 samples.

\subsection{Confidence Stratification}
\label{sec:alg-confidence}

\begin{algorithm}[H]
\caption{Confidence Stratification}
\label{alg:confidence}
\DontPrintSemicolon
\KwIn{Validity indicators $\epsilon_1, \epsilon_2, \epsilon_3, \epsilon_4$}
\KwOut{Confidence tier $\in \{\text{high, moderate, low}\}$, flags}
\BlankLine
flags $\gets []$ \;
\lIf{$\epsilon_1 > 0.5$}{flags.append(``PERTURBATIVE\_BREAKDOWN'')}
\lIf{$\epsilon_2 > 0.3$}{flags.append(``MOMENT\_CLOSURE\_POOR'')}
\lIf{$\epsilon_3 > 0.15$}{flags.append(``TWO\_LOOP\_DIVERGENCE'')}
\lIf{$\epsilon_4 > 0.1$}{flags.append(``NYSTROM\_ERROR\_HIGH'')}
\BlankLine
$n_{\mathrm{violations}} \gets |\text{flags}|$ \;
\lIf{$n_{\mathrm{violations}} = 0$}{tier $\gets$ ``high''}
\lElseIf{$n_{\mathrm{violations}} = 1$}{tier $\gets$ ``moderate''}
\lElse{tier $\gets$ ``low''}
\Return tier, flags
\end{algorithm}

% =============================================================================
\section{Computational Complexity Analysis}
\label{sec:complexity}

We analyze the end-to-end computational cost of the phase diagram pipeline.

\paragraph{Stage 1: Calibration (amortized).} NTK computation at $W = 5$ widths with $S = 10$ seeds: $\cO(WS \cdot n^2 L N_{\max}^2) = \cO(50 \cdot 200^2 \cdot 4 \cdot 1024^2) \approx 8.4 \times 10^{12}$ FLOPs. At $\sim 10^{10}$ FLOP/s on a modern CPU: $\sim 14$ minutes. This is computed \emph{once} per architecture family.

\paragraph{Stage 2: Per-grid-point.} For each $(\eta, N)$ pair:
\begin{itemize}[leftmargin=*,itemsep=0pt]
\item ODE integration: $\cO(K_{\mathrm{steps}} \cdot n^3) \approx 500 \cdot 200^3 = 4 \times 10^9$ FLOPs $\approx 0.4$ seconds.
\item Lyapunov exponent: $\cO(1000 \cdot n^3) \approx 8 \times 10^9$ FLOPs $\approx 0.8$ seconds.
\item Total per point: $\sim 1.2$ seconds.
\end{itemize}

\paragraph{Stage 3: Boundary detection.} Each bisection requires $\sim 25$ growth rate evaluations: $\sim 20$ seconds per boundary point. Continuation with $K = 50$ steps: $\sim 50 \times 60 = 3000$ growth rate evaluations $\approx 40$ minutes per boundary curve.

\paragraph{Full pipeline.} For a grid of $|\eta| \times |N| = 20 \times 5 = 100$ points: Stage 2 takes $\sim 2$ minutes. Boundary detection for $\sim 5$ boundary curves: $\sim 3$ hours. Calibration: $\sim 15$ minutes. \textbf{Total system prediction time: $\sim 2$--$4$ hours on a single CPU.}

\paragraph{Comparison to full training.} Ground truth requires training $\sim 4000$ networks (Section~\ref{sec:ground-truth}) at $\sim 15$ minutes each: $\sim 1000$ CPU-hours. The phase diagram system provides predictions at $\sim 1/250$ the cost. More importantly, the system produces predictions \emph{without any training runs}, given only the architecture and dataset.

\paragraph{Bottleneck analysis.} The bottleneck is the Lyapunov exponent computation within pseudo-arclength continuation ($\sim 90\%$ of total time). This scales as $\cO(n^3)$ per evaluation. At $n_{\mathrm{probe}} = 200$, this is tractable; at $n = 2000$, the $1000 \times$ increase would require Nystr\"om approximation or matrix-free methods to remain feasible.

% =============================================================================
\section{Experimental Design}
\label{sec:experiments}

We present the complete experimental plan for validating the framework. This section specifies protocols, metrics, and baselines; numerical results are deferred to the implementation stage.

\subsection{Ground Truth Protocol}
\label{sec:ground-truth}

\paragraph{Operational regime definition.} A configuration (architecture, width, learning rate, dataset) is classified using a \emph{multi-criteria conjunction} of four indicators:

\begin{center}
\begin{tabular}{@{}llll@{}}
\toprule
\# & Indicator & Lazy criterion & Rich criterion \\
\midrule
1 & KADR (Def.~\ref{def:KADR}) & $\bar{\rho}_{\mathrm{avg}} < \tau_1$ & $\bar{\rho}_{\mathrm{avg}} \ge \tau_1$ \\
2 & CKA (init vs.\ final) & $\CKA > 0.90$ & $\CKA \le 0.90$ \\
3 & Weight displacement & $\norm{W_f - W_0}_F/\norm{W_0}_F < 0.10$ & $\ge 0.10$ \\
4 & Linear probe gap & $|\text{Acc}_{\mathrm{net}} - \text{Acc}_{\mathrm{probe}}| < 0.05$ & $\ge 0.05$ \\
\bottomrule
\end{tabular}
\end{center}

A configuration is labeled \textbf{rich} if $\ge 3$ of 4 indicators agree; \textbf{lazy} if $\ge 3$ indicate lazy; \textbf{ambiguous} if 2/2 split (excluded from AUC but reported separately). The threshold $\tau_1$ is set adaptively per architecture family using 10 known-lazy and 10 known-rich calibration configurations.

\paragraph{Non-circularity argument.} The system predicts via perturbative kernel evolution ODEs (no training). Ground truth uses actual gradient descent training with empirical NTK recomputation. These are methodologically independent: the system never sees training trajectories; ground truth never uses perturbation theory. The shared use of KADR as a measurement quantity is a \emph{definitional} choice (measuring the same physical quantity via different computational paths), not a \emph{computational} circularity. Indicators 2--4 (CKA, weight displacement, linear probe gap) are entirely independent of NTK and provide a check on this definitional coupling.

\paragraph{Seed counts.} Boundary-adjacent configurations ($|\log(\eta/\eta^*)| < 0.3$): 20 seeds. Interior configurations: 5 seeds. Calibration runs: 10 seeds. Total: $\sim 4000$ training runs.

\paragraph{Adaptive boundary refinement.} Phase~1 (coarse): 5 seeds on 8-point log-uniform grid per axis. Phase~2: fit sigmoid to identify $\eta^*$ estimate. Phase~3: add 5 grid points in $[\eta^*/3, 3\eta^*]$. Phase~4: augment to 20 seeds near $\eta^*$. Phase~5: re-fit on refined grid.

\subsection{Evaluation Metrics}
\label{sec:metrics}

\begin{enumerate}[leftmargin=*]
\item \textbf{Classification AUC.} Area under ROC curve for lazy-vs-rich binary classification. DeLong confidence intervals. Target: AUC $> 0.90$.

\item \textbf{Boundary localization error.} $\mathrm{RelErr} = |\log(\eta^*_{\mathrm{pred}}/\eta^*_{\mathrm{true}})| / |\log(\eta^*_{\mathrm{true}}/\eta_{\min})|$, normalizing by dynamic range. Target: $< 15\%$ relative error for $> 80\%$ of configurations.

\item \textbf{Confidence-stratified accuracy.} Report accuracy separately for high, moderate, and low confidence predictions. Target: $> 95\%$ accuracy for high-confidence, $> 80\%$ for moderate.

\item \textbf{Scaling exponent prediction error.} For the finite-size scaling prediction $\Delta\eta/\eta^* \sim N^{-\beta}$: $|\beta_{\mathrm{pred}} - \beta_{\mathrm{meas}}| / \beta_{\mathrm{meas}}$. Report with 95\% CI.

\item \textbf{UQ coverage.} Fraction of prediction failures flagged as low-confidence. Target: $> 90\%$.
\end{enumerate}

\subsection{Baselines}
\label{sec:baselines}

\begin{enumerate}[leftmargin=*]
\item \textbf{Baseline A: Infinite-width NTK.} Predicts ``always lazy.'' Accuracy equals the lazy base rate ($\sim 50\%$ on balanced test set). Implemented via Neural Tangents library.

\item \textbf{Baseline B: $\muP$ transfer.} Predicts regime based on $\gamma = \eta N^{-(1-a-b)} \gtrless C_{\mathrm{fit}}$ with threshold $C_{\mathrm{fit}}$ tuned on a held-out calibration set. Captures the correct scaling but not architecture/data dependence.

\item \textbf{Baseline C: Mean-field (two-layer only).} For two-layer networks, uses the Chizat--Bach condition $\eta\sqrt{N} > C$ \citep{chizat2019lazy}. Inapplicable to deeper architectures.

\item \textbf{Baseline D: Empirical early-KADR.} Train at the target width for 200 steps, measure KADR, extrapolate to 2000 steps. Cheap empirical approach that could outperform the ODE machinery if early dynamics are predictive.

\item \textbf{Baseline E: Width interpolation (strengthened).} Given the \emph{same} 5 calibration widths and 10 seeds as the system, fit $\KADR(N) = a + b/N$ by linear regression and predict at target width. This is a strong baseline that uses actual training data; the system must demonstrate that ODE-based prediction adds value beyond interpolation.
\end{enumerate}

\paragraph{Fair comparison protocol.} All baselines use the same calibration data (5 widths, 10 seeds). Baseline E is given identical training budget. Statistical comparison via McNemar's test (paired, Bonferroni-corrected to $\alpha' = 0.01$ for 5 baselines).

\subsection{Ablation Studies}
\label{sec:ablations}

Four mandatory ablations isolate the contribution of each component:

\begin{enumerate}[leftmargin=*]
\item \textbf{$1/N$ corrections ON vs.\ OFF.} Replace $\Theta^{(0)} + \Theta^{(1)}/N$ with $\Theta^{(0)}$ alone. Measures whether finite-width corrections improve prediction beyond the infinite-width kernel.

\item \textbf{Frozen-$H$ vs.\ periodically recomputed $H$.} For 10\% of configurations, compare the frozen-$H$ ODE against periodic $H$ recomputation every 100 gradient steps. Validates Lemma~\ref{lem:structural-stability}.

\item \textbf{Calibration widths $\{64, 128, 256\}$ vs.\ $\{64, 128, 256, 512, 1024\}$.} Quantifies the value of additional calibration points and tests sensitivity to the number of widths.

\item \textbf{Moment closure sensitivity.} Perturb $\kappa_4 \to \delta \cdot \kappa_4^{(\mathrm{est})}$ for $\delta \in \{-0.5, 0, 0.5\}$ and measure the resulting shift in $\gamma_c$. Reports the sensitivity coefficient $\chi_{\kappa_4}$.
\end{enumerate}

\subsection{Retrodiction Validation}
\label{sec:retrodictions}

The system must recover the following known results:

\begin{enumerate}[leftmargin=*]
\item \textbf{Chizat--Bach limit.} For two-layer ReLU networks, the lazy-to-rich boundary scales as $\eta \propto 1/\sqrt{N}$ under standard parameterization \citep{chizat2019lazy}. Our system should recover this scaling.

\item \textbf{Saxe et al.\ linear dynamics.} For deep linear networks, training dynamics are exactly solvable \citep{saxe2014exact}. Our kernel ODE should reproduce the known loss curves and kernel evolution.

\item \textbf{$\muP$ scaling.} Under maximal update parameterization, the phase boundary should be independent of $N$ when expressed in terms of $\gamma$ \citep{yang2021tensor}.
\end{enumerate}

\subsection{Novel Predictions}
\label{sec:novel-predictions}

\begin{enumerate}[leftmargin=*]
\item \textbf{Phase boundary sharpness scaling.} The transition width $\Delta\eta$ (interval where KADR goes from 20\% to 80\% of maximum) scales as $\Delta\eta/\eta^* \sim N^{-\beta}$ with $\beta \in [0.3, 0.7]$ for MLPs. Null hypothesis: $\beta = 0$ (no sharpening). Test via OLS regression of $\log(\Delta\eta/\eta^*)$ on $\log(N)$.

\item \textbf{Architecture-dependent critical exponents.} If the exponent $\beta$ differs between MLP, ConvNet, and ResNet families (with $p < 0.01$ by $F$-test on the regression slopes), this provides evidence for distinct \emph{universality classes} in training dynamics---connecting neural network theory to statistical mechanics.

\item \textbf{Dataset complexity determines boundary.} $\eta^* \sim d_{\mathrm{eff}}^{-\alpha}$ where $d_{\mathrm{eff}}$ is the effective rank of the target function in the kernel eigenbasis.

\item \textbf{Depth-dependent boundary shift.} $\eta^*(L) \sim L^{-\delta}$ with $\delta > 0$, predicting that deeper networks transition to the rich regime at smaller learning rates.
\end{enumerate}

\subsection{Definition Concordance Experiment}
\label{sec:concordance}

Three definitions of the phase boundary coexist: (1) Lyapunov exponent zero-crossing (Theorem~\ref{thm:stability-boundary}), (2) steepest-gradient of KADR (Definition~\ref{def:operational-boundary}), (3) sigmoid inflection of KADR vs.\ $\eta$. On 3 calibration architectures (MLP-2, MLP-3, MLP-4 on MNIST, width 256), we compute all three definitions and verify agreement within measurement uncertainty (bootstrap 95\% CI). Disagreement $> 10\%$ relative in $\gamma_c$ triggers investigation.

This experiment validates Theorem~\ref{thm:stability-boundary}(iii) and justifies the use of the steepest-gradient locus as the primary operational definition.

\subsection{Frozen-$H$ Validation}
\label{sec:frozen-H-validation}

For a representative configuration near the predicted boundary in each architecture family:

\begin{enumerate}[leftmargin=*]
\item \textbf{Frozen-$H$ trajectory:} Solve ODE with $H = H(0)$ (Algorithm~\ref{alg:ode}).
\item \textbf{Periodically-recomputed-$H$:} Re-evaluate $H(\theta(t))$ every 100 gradient steps by integrating the full parameter ODE and recomputing $H$ from updated parameters.
\item \textbf{Ground truth:} Full training with empirical NTK tracking.
\end{enumerate}

Compare predicted $\gamma_c$ from methods 1 and 2. If boundary shift $> 15\%$, the frozen-$H$ approximation is unreliable near the boundary, and Lemma~\ref{lem:structural-stability} requires tighter constants.

% =============================================================================
\section{Connections to Existing Theory}
\label{sec:connections}

\paragraph{NTK theory at infinite width.} In the limit $N \to \infty$ (equivalently $\gamma \to 0$ under SP), our framework recovers the standard NTK theory: $\Theta^{(1)} \to 0$, the kernel is frozen, and training dynamics are exactly those of kernel regression with $\Theta^{(0)}$. The perturbative correction provides the leading finite-width departure from this limit.

\paragraph{Mean-field theory.} Mean-field theory operates in the complementary regime: $\gamma \to \infty$ (or equivalently, the rich regime deep beyond the phase boundary). Our framework describes the \emph{onset} of the mean-field regime, providing the critical coupling $\gamma_c$ at which the NTK description breaks down. The two frameworks are complementary rather than competing: ours predicts \emph{when} feature learning begins; mean-field theory describes what happens \emph{after} it begins.

\paragraph{$\muP$ and hyperparameter transfer.} The $\muP$ framework \citep{yang2021tensor} identifies the correct scaling exponents $(a,b)$ for hyperparameter transfer. Our dimensionless coupling $\gamma = \eta N^{-(1-a-b)}$ provides a mechanistic explanation: $\muP$ works because it holds $\gamma$ fixed while scaling $N$, keeping the system at the same point in the phase diagram. The phase boundary $\gamma_c$ is the quantitative prediction that $\muP$ alone does not provide.

\paragraph{Edge of stability.} The edge-of-stability phenomenon \citep{cohen2021gradient} arises from discrete gradient descent (not gradient flow) and involves a Hopf-type instability where the loss sharpness hovers near $2/\eta$. In our framework, this corresponds to the imaginary part of the critical eigenvalue becoming nonzero---a Hopf bifurcation of the discrete-time map. Our continuous-time (gradient flow) analysis captures the transcritical instability but not the oscillatory dynamics. Extending to discrete updates (replacing A2) would connect the two phenomena; this is deferred to future work.

\paragraph{Statistical mechanics of learning.} Our finite-size scaling prediction ($\Delta\eta/\eta^* \sim N^{-\beta}$) is the neural network analog of finite-size scaling in statistical mechanics \citep{fisher1972scaling}. The width $N$ plays the role of system size, $\gamma$ the role of temperature or coupling constant, and $\beta$ the critical exponent. If different architecture families (MLP, ConvNet, ResNet) have different $\beta$ values, they belong to different \emph{universality classes}---a concept that would import the powerful classification theory of critical phenomena into deep learning.

% =============================================================================
\section{Scope and Limitations}
\label{sec:limitations}

We provide an honest accounting of the framework's limitations.

\paragraph{Transformer exclusion.} Softmax attention creates a fundamentally different NTK structure (data-dependent, with rank deficiency from the softmax normalization) that is not captured by our perturbative expansion. Extending the framework to transformers requires a separate analysis of the attention kernel and its finite-width corrections, which is beyond the scope of this work.

\paragraph{Moment closure quality.} The Gaussian moment closure (Assumption~\ref{ass:moment-closure}) sets $\kappa_4 = 0$, which is well-justified for GELU (nearly Gaussian pre-activations) but introduces $\cO(\kappa_4/N)$ systematic bias for ReLU (half-Gaussian after one layer). The sensitivity coefficient $\chi_{\kappa_4}$ quantifies this; if large, the phase boundary prediction is dominated by an uncontrolled approximation. We mitigate by (1) using GELU as the primary activation and (2) empirically measuring $\kappa_4/\kappa_2^2$ at initialization.

\paragraph{Perturbative regime restrictions.} Near the phase boundary, the perturbative expansion is \emph{least} reliable---precisely where predictions are most valuable. The validity indicator $\epsilon_1$ flags when the expansion ratio exceeds 0.5, but the system cannot guarantee accuracy at the boundary itself. The frozen-$H$ ODE correctly identifies the \emph{location} of the boundary (as the onset of instability) even when quantitative trajectory predictions degrade.

\paragraph{Exponential validity bound.} Theorem~\ref{thm:validity} provides a Gr\"onwall-type bound that is exponential in $T$, rendering the theoretical minimum width cosmologically large for practical training durations. The bound is useful qualitatively (indicating degradation with $\gamma, T$, conditioning) but the quantitative diagnostic must rely on the computable indicator $\epsilon_1$.

\paragraph{Dataset subsampling.} All NTK computations use a probe set of $n_{\mathrm{probe}} = 200$ samples. This is sufficient for KADR estimation but may miss fine-grained spectral structure of the NTK on the full dataset. For datasets with complex cluster structure, the probe set may not be representative.

\paragraph{Gradient flow assumption.} Continuous-time dynamics (Assumption~\ref{ass:gradient-flow}) introduce $\cO(\eta^2)$ errors relative to discrete gradient descent. At large learning rates ($\eta > 0.1$), edge-of-stability effects and oscillatory dynamics are not captured. We flag predictions at $\eta > 0.1$ as ``moderate confidence'' due to this approximation.

\paragraph{Depth limitations.} The perturbative expansion has effective parameter $L/N$. For $L = \cO(N)$, the expansion breaks down regardless of individual layer widths. Our experiments are restricted to $L \le 6$, ensuring $L/N \le 6/64 \approx 0.1$ at the minimum width.

\paragraph{Regression conditioning.} The WLS regression for extracting $\Theta^{(1)}$ (Algorithm~\ref{alg:correction}) has condition number $\cO(N_{\max})$ after centering, which is manageable ($\sim 10^3$ at $N_{\max} = 1024$). Without centering, the condition number is $\cO(N_{\max}^2) \sim 10^6$, which can produce unreliable slope estimates. We report the condition number of $A^\top WA$ and flag if $> 10^4$.

\paragraph{No mini-batch SGD.} All experiments use full-batch gradient descent. Mini-batch SGD introduces gradient noise that can qualitatively change the phase diagram (noise-induced transitions, implicit regularization). Extending to stochastic dynamics is a natural but non-trivial extension.

% =============================================================================
\section{Conclusion}
\label{sec:conclusion}

We have developed a systematic framework for computing finite-width phase diagrams that predict the lazy-to-rich transition in neural network training. The key theoretical contributions are: (1)~the $H_{ijk}$ tensor decompositions for convolutional and residual architectures, providing the first analytical handle on finite-width kernel evolution for weight-shared networks; (2)~the Lyapunov stability formulation of the phase boundary, resolving a fundamental issue with prior steady-state bifurcation analyses; and (3)~the structural stability lemma justifying the frozen-$H$ approximation that makes computation tractable.

The framework enables three capabilities that did not previously exist: \emph{prediction} of training regime from architecture and dataset alone (without training), \emph{quantification} of the sharpness and location of phase transitions as a function of width, and \emph{diagnosis} of when perturbative predictions are reliable through self-diagnostic validity indicators.

The most impactful potential finding is the existence of architecture-dependent universality classes in training dynamics. If the finite-size scaling exponent $\beta$ differs measurably across MLP, ConvNet, and ResNet families, this would establish a deep connection between neural network training and the theory of critical phenomena in statistical mechanics---providing a principled classification of architectures based on their dynamical phase structure.

\paragraph{Future directions.} Natural extensions include: (i)~discrete gradient descent dynamics (connecting to edge-of-stability theory), (ii)~mini-batch SGD (adding noise-induced phase transitions), (iii)~transformer architectures (requiring attention-kernel analysis), (iv)~higher-loop corrections (enabling predictions at narrower widths), and (v)~cross-entropy loss (modifying the residual evolution equation). Each extension addresses a specific limitation of the current framework while building on its mathematical infrastructure.

% =============================================================================
\appendix

\section{Proof Details}
\label{app:proofs}

\subsection{Detailed Proof of Theorem~\ref{thm:stability-boundary}, Part~(ii): Monotonicity}

We provide the detailed argument for the monotonicity of $\Lambda_T(\gamma)$ in $\gamma$.

\begin{proof}
Write the variational system \eqref{eq:variational} in matrix form as
\begin{equation}
\frac{\dd}{\dd t}\begin{pmatrix} \delta\Theta_{\mathrm{vec}} \\ \delta r \end{pmatrix} = A(t;\gamma) \begin{pmatrix} \delta\Theta_{\mathrm{vec}} \\ \delta r \end{pmatrix},
\end{equation}
where $\delta\Theta_{\mathrm{vec}}$ is the vectorization of the upper triangle of $\delta\Theta$, and
\begin{equation}
A(t;\gamma) = \begin{pmatrix} 0 & -\eta\gamma\, \tilde{H}(t) \\ -\eta\gamma\, B(t) & -\eta\,\Theta^{(0)} \end{pmatrix}.
\end{equation}
Here $\tilde{H}(t)$ encodes the $H$-tensor contraction with $r_0(t)$: $[\tilde{H}(t)]_{\mathrm{vec}(ij),k} = r_{0,k}(t) H_{ijk}$, and $B(t)$ encodes the kernel-residual coupling: $[B(t)]_{i,\mathrm{vec}(jk)} = r_{0,j}(t)\,\delta_{ik}$ (schematically).

For $\gamma_2 > \gamma_1 > 0$, define $\Delta A = A(t;\gamma_2) - A(t;\gamma_1) = (\gamma_2 - \gamma_1) A_1(t)$ where
\begin{equation}
A_1(t) = \begin{pmatrix} 0 & -\eta\,\tilde{H}(t) \\ -\eta\,B(t) & 0 \end{pmatrix}.
\end{equation}
The comparison result for finite-time Lyapunov exponents states that if $\Phi_2(T,0)$ and $\Phi_1(T,0)$ are the state-transition matrices, then
\begin{equation}
\ln\norm{\Phi_2(T,0)}_{\mathrm{op}} \ge \ln\norm{\Phi_1(T,0)}_{\mathrm{op}} - \epsilon,
\end{equation}
where $\epsilon$ accounts for the non-commutativity of $A(t;\gamma)$ at different times. In the present case, the perturbation $\Delta A$ enhances the off-diagonal coupling (which drives perturbation growth) without affecting the diagonal dissipation ($-\eta\,\Theta^{(0)}$). A more refined argument uses the Floquet--Lyapunov theory for non-autonomous systems.

Consider the quadratic functional $V(t) = \norm{(\delta\Theta(t), \delta r(t))}^2$. Its time derivative is
\begin{equation}
\dot{V} = 2\ip{(\delta\Theta, \delta r)}{A(t;\gamma)(\delta\Theta, \delta r)}.
\end{equation}
The off-diagonal blocks of $A$ couple $\delta\Theta$ and $\delta r$, and increasing $\gamma$ increases the magnitude of these coupling terms. Since the coupling enables energy transfer from the dissipated $\delta r$ direction (which has growth rate $-\eta\lambda_{\min}(\Theta^{(0)}) < 0$) into the $\delta\Theta$ direction (which has zero intrinsic growth rate but can grow through the feedback loop), increasing $\gamma$ strictly increases the maximum growth rate.

Formally, for any unit initial condition $v_0$, the growth functional $G(\gamma, v_0) = (1/T)\ln\norm{\Phi_\gamma(T,0)v_0}$ satisfies $\pp G/\pp\gamma > 0$ whenever the trajectory achieves net growth through the feedback mechanism. Since $\Lambda_T = \sup_{v_0} G$, and the supremum is achieved by the most unstable direction (which benefits from increased coupling), monotonicity follows.
\end{proof}

\subsection{Proof of Lemma~\ref{lem:spatial-averaging}: Berry--Esseen Refinement}

\begin{proof}
Define $Z_\alpha = \Delta H^{(\mathrm{dense})}_{\alpha\beta}(x_i, x_j, x_k)$ for fixed $i,j,k$ and varying $(\alpha,\beta) \in \cS^2 \setminus \{(\alpha,\alpha)\}$. Under translation covariance, $\E[Z_\alpha] = 0$ and
\begin{equation}
\Cov(Z_\alpha, Z_{\alpha'}) = \sigma^2 \cdot \rho(|\alpha - \alpha'|),
\end{equation}
where $\rho(r) = e^{-r/\xi}$ is the spatial correlation function with correlation length $\xi$. The spatial correction is $\Delta H^{(\mathrm{spatial})} = \sum_{\alpha \neq \beta} Z_{\alpha\beta}$.

The variance of this sum is
\begin{align}
\Var\!\left[\sum_{\alpha \neq \beta} Z_{\alpha\beta}\right]
&= \sum_{\alpha \neq \beta} \sum_{\alpha' \neq \beta'} \Cov(Z_{\alpha\beta}, Z_{\alpha'\beta'}) \\
&\le |\cS|^2 \cdot \sigma^2 \cdot \sum_{r=-|\cS|}^{|\cS|} |\cS| \cdot \rho(|r|) \\
&= |\cS|^3 \cdot \sigma^2 \cdot 2\xi \cdot (1 + \cO(e^{-|\cS|/\xi})).
\end{align}
Therefore $\norm{\Delta H^{(\mathrm{spatial})}}_F = \cO(|\cS|^{3/2} \cdot \sigma \cdot \sqrt{\xi})$.

Dividing by $\norm{\bar{H}^{(\mathrm{dense})}}_F = \cO(|\cS| \cdot \sigma)$ (from the $|\cS|$ diagonal terms, each of magnitude $\sigma$):
\begin{equation}
\frac{\norm{\Delta H^{(\mathrm{spatial})}}_F}{\norm{\bar{H}^{(\mathrm{dense})}}_F} = \cO\!\left(\frac{|\cS|^{3/2} \sqrt{\xi}}{|\cS|}\right) = \cO\!\left(\frac{\sqrt{\xi \cdot |\cS|}}{|\cS|}\right) = \cO\!\left(\frac{\sqrt{\xi}}{\sqrt{|\cS|}}\right) = \cO\!\left(\frac{\xi}{\sqrt{|\cS|}}\right).
\end{equation}
The last step uses $\sqrt{\xi}/\sqrt{|\cS|} \le \xi/\sqrt{|\cS|}$ for $\xi \ge 1$ (correlation length at least one grid spacing). For $\xi < 1$, the bound is even stronger.
\end{proof}

\subsection{Detailed Constants in Theorem~\ref{thm:validity}}

The constants $C_1$ and $C_2$ in the validity bound \eqref{eq:validity-bound} are:
\begin{align}
C_1 &= L \cdot C_\sigma^3 \cdot n \cdot \max_\ell \frac{\norm{K^{(\ell)}}_{\mathrm{op}}^2}{\lambda_{\min}(\Theta^{(0)})}, \\
C_2 &= L \cdot C_\sigma \cdot \left(1 + \frac{C_\sigma}{\lambda_{\min}(\Theta^{(0)})}\right),
\end{align}
where $C_\sigma = \max(C_0, C_1, C_2, C_3)$ bounds the activation derivatives (Assumption~\ref{ass:smooth}), $K^{(\ell)}$ is the layer-$\ell$ kernel matrix, and $n$ is the dataset size. These can be computed numerically for any specific architecture and dataset.

For GELU: $C_0 \le 1.13$, $C_1 = 1$ (Lipschitz constant), $C_2 \le 0.82$, $C_3 \le 0.48$, so $C_\sigma = 1.13$.

\section{Algorithm Implementation Details}
\label{app:implementation}

\subsection{Numerical Stability of the ODE Solver}

The coupled system \eqref{eq:coupled-system} can be stiff when the spectral ratio $\lambda_{\max}(\Theta^{(0)})/\lambda_{\min}(\Theta^{(0)})$ is large (i.e., the NTK is ill-conditioned). We use the following heuristics:

\begin{itemize}[leftmargin=*]
\item \textbf{Stiffness detection:} Compute the spectral radius $\eta\,\lambda_{\max}(\Theta^{(0)})$ and compare to $1/T$. If the stiffness ratio $\eta\,\lambda_{\max}\,T > 1000$, switch from explicit RK45 to implicit BDF.
\item \textbf{Tolerance selection:} Use relative tolerance $\epsilon_{\mathrm{rtol}} = 10^{-6}$ and absolute tolerance $\epsilon_{\mathrm{atol}} = 10^{-8}$. The relative tolerance controls the fast modes (large eigenvalues) and the absolute tolerance prevents the slow modes from being lost in roundoff.
\item \textbf{Maximum step size:} Set $\Delta t_{\max} = T/10$ to prevent the adaptive integrator from taking steps that miss transient dynamics near the phase boundary.
\end{itemize}

\subsection{Lyapunov Exponent Convergence}

The QR-based algorithm (Algorithm~\ref{alg:growth-rate}) requires choosing the number of reorthogonalization intervals $M$ and the number of exponents $k$. We use:

\begin{itemize}[leftmargin=*]
\item $M = 100$: sufficient for convergence of the dominant Lyapunov exponent to within $\sim 1\%$ for training horizons $T = 2000$.
\item $k = 10$: captures the dominant growth directions. In practice, only $\Lambda_1$ (the largest exponent) is needed for phase boundary detection; $\Lambda_2, \ldots, \Lambda_k$ provide additional diagnostic information about the instability structure.
\item Convergence check: compute $\Lambda_1$ at $M$ and $2M$ intervals; if they differ by $> 5\%$, double $M$.
\end{itemize}

\subsection{Continuation Step-Size Adaptation}

The pseudo-arclength continuation (Algorithm~\ref{alg:continuation}) uses adaptive step-size control based on three criteria:

\begin{enumerate}[leftmargin=*]
\item \textbf{Curvature:} If the angle between successive tangent vectors exceeds $60°$ ($\cos\alpha < 0.5$), halve the step size. If the angle is less than $18°$ ($\cos\alpha > 0.95$), increase by 50\%.
\item \textbf{Newton convergence:} If the corrector requires more than 10 iterations, halve the step size for the next step.
\item \textbf{Fold-point detection:} If $\det(J_{\mathrm{aug}}) < 10^{-14}$, a fold point (turning point) is nearby. Switch to a specialized fold-point handler that detects the fold precisely via an augmented system.
\end{enumerate}

Bounds: $\Delta s_{\min} = 10^{-8}$, $\Delta s_{\max} = 0.5$ (in normalized parameter space).

\subsection{H-Tensor Computation}

For MLP architectures where the analytical form is known \citep{dyer2020asymptotics}, we compute $H_{ijk}$ directly from the layer-wise kernel matrices and their derivatives. For convolutional and residual architectures, we use a hybrid approach:

\begin{enumerate}[leftmargin=*]
\item \textbf{Analytical computation} (Theorems~\ref{thm:H-conv}, \ref{thm:H-skip}) when the architecture matches the theorem assumptions exactly.
\item \textbf{Numerical validation:} At small widths ($N \in \{8, 16, 32\}$), compare the analytical $H$ against a brute-force finite-difference computation: $H_{ijk} \approx (\Theta_{ij}(\theta + \epsilon e_k) - \Theta_{ij}(\theta - \epsilon e_k)) / (2\epsilon)$, aggregated over parameter directions. Agreement within 5\% relative error validates the analytical formula.
\item \textbf{Fallback:} If the analytical formula is not available or fails validation, use the empirical $H$ estimated from multi-width kernel regression.
\end{enumerate}

\section{Assumption Verification Protocols}
\label{app:assumptions}

\begin{center}
\begin{tabular}{@{}p{2.5cm}p{4.5cm}p{5cm}@{}}
\toprule
Assumption & Verification method & Action if violated \\
\midrule
A1 (Gaussian init) & By construction (Kaiming/He) & Use exact init; no action needed \\
A2 (Gradient flow) & Compare $\eta$ and $\eta/10$ predictions for 10\% of configs & Flag $\eta > 0.1$ as moderate confidence \\
A3 (Squared loss) & By construction (MSE) & No action needed \\
A4 (Uniform width) & Compute $\max_\ell N_\ell / \min_\ell N_\ell$ & Use $N_{\mathrm{eff}} = \min_\ell N_\ell$; report if ratio $> 2$ \\
A5 ($1/N$ convergence) & Two-loop check $\epsilon_3$ & Restrict claims to widths with $\epsilon_3 < 0.15$ \\
A6 (Moment closure) & MC estimate of $\kappa_4/\kappa_2^2$ at init & Flag if $> 0.3$; for ReLU, acknowledge bias \\
A7 (Layer independence) & By construction (no BatchNorm) & Exclude normalization layers \\
A8$^*$ ($C^3$ activation) & By construction (GELU primary) & ReLU: Conjecture~\ref{conj:relu} only \\
A9 (Data regularity) & Compute $\kappa(X^\top X/d)$ & Remove near-duplicates; subsample if needed \\
A10 (Finite horizon) & By construction ($T = 2000$) & Monitor $\epsilon_1(t)$ for secular growth \\
\bottomrule
\end{tabular}
\end{center}

\section{Additional Known Results Used}
\label{app:known-results}

For completeness, we state two known lemmas that are used in the proofs.

\begin{lemma}[NTK positive definiteness at finite width {\citep{du2019gradient,nguyen2021tight}}]
\label{lem:ntk-pd}
Under Assumptions~\ref{ass:gaussian}, \ref{ass:smooth}, \ref{ass:data}, for any $L$-layer network with width $N \ge N_0(n, L, \sigma)$, the empirical NTK at initialization is positive definite with probability $\ge 1 - \delta$:
\begin{equation}
\lambda_{\min}(\Theta(\theta_0)) \ge \tfrac{1}{2}\lambda_{\min}(\Theta^{(0)}) > 0,
\end{equation}
where $N_0 = \cO(n^2 L^2 \log(n/\delta))$.
\end{lemma}

\begin{lemma}[Nystr\"om error bound {\citep{gittens2016revisiting}}]
\label{lem:nystrom-error}
For a positive semidefinite matrix $\Theta^{(0)} \in \R^{n \times n}$ with eigenvalues $\lambda_1 \ge \cdots \ge \lambda_n \ge 0$, the rank-$m$ Nystr\"om approximation with $k$-means landmarks satisfies:
\begin{equation}
\norm{\Theta^{(0)} - \Theta^{(0)}_m}_2 \le (1 + \epsilon_{\mathrm{kmeans}})\,\lambda_{m+1}.
\end{equation}
The induced phase boundary shift satisfies:
\begin{equation}
|\gamma_c(\Theta^{(0)}_m) - \gamma_c(\Theta^{(0)})| \le C \cdot \frac{\lambda_{m+1}}{\lambda_{\min}(\Theta^{(0)})} \cdot \gamma_c(\Theta^{(0)}).
\end{equation}
\end{lemma}

% =============================================================================
% References
% =============================================================================

\begin{thebibliography}{30}

\bibitem[Jacot et~al.(2018)]{jacot2018neural}
A.~Jacot, F.~Gabriel, and C.~Hongler.
\newblock Neural tangent kernel: Convergence and generalization in neural networks.
\newblock In \emph{Advances in Neural Information Processing Systems (NeurIPS)}, 2018.

\bibitem[Dyer \& Gur-Ari(2020)]{dyer2020asymptotics}
E.~Dyer and G.~Gur-Ari.
\newblock Asymptotics of wide networks from {F}eynman diagrams.
\newblock In \emph{International Conference on Learning Representations (ICLR)}, 2020.

\bibitem[Huang \& Yau(2020)]{huang2020dynamics}
J.~Huang and H.-T. Yau.
\newblock Dynamics of deep neural networks and neural tangent hierarchy.
\newblock In \emph{International Conference on Machine Learning (ICML)}, 2020.

\bibitem[Mei et~al.(2018)]{mei2018mean}
S.~Mei, A.~Montanari, and P.-M.~Nguyen.
\newblock A mean field view of the landscape of two-layer neural networks.
\newblock \emph{Proceedings of the National Academy of Sciences}, 115(33):E7665--E7671, 2018.

\bibitem[Yang \& Hu(2021)]{yang2021tensor}
G.~Yang and E.~J.~Hu.
\newblock Tensor programs {IV}: Feature learning in infinite-width neural networks.
\newblock In \emph{International Conference on Machine Learning (ICML)}, 2021.

\bibitem[Chizat \& Bach(2019)]{chizat2019lazy}
L.~Chizat and F.~Bach.
\newblock A note on lazy training in supervised differentiable programming.
\newblock \emph{arXiv preprint arXiv:1812.07956}, 2019.

\bibitem[Rotskoff \& Vanden-Eijnden(2018)]{rotskoff2018parameters}
G.~M.~Rotskoff and E.~Vanden-Eijnden.
\newblock Parameters as interacting particles: long time convergence and asymptotic error scaling of neural networks.
\newblock In \emph{Advances in Neural Information Processing Systems (NeurIPS)}, 2018.

\bibitem[Chizat \& Bach(2018)]{chizat2018global}
L.~Chizat and F.~Bach.
\newblock On the global convergence of gradient descent for over-parameterized models using optimal transport.
\newblock In \emph{Advances in Neural Information Processing Systems (NeurIPS)}, 2018.

\bibitem[Cohen et~al.(2021)]{cohen2021gradient}
J.~M.~Cohen, S.~Kaur, Y.~Li, J.~Z.~Kolter, and A.~Talwalkar.
\newblock Gradient descent on neural networks typically occurs at the edge of stability.
\newblock In \emph{International Conference on Learning Representations (ICLR)}, 2021.

\bibitem[Lewkowycz et~al.(2020)]{lewkowycz2020large}
A.~Lewkowycz, Y.~Bahri, E.~Dyer, J.~Sohl-Dickstein, and G.~Gur-Ari.
\newblock The large learning rate phase of neural network training.
\newblock \emph{arXiv preprint arXiv:2002.10505}, 2020.

\bibitem[Woodworth et~al.(2020)]{woodworth2020kernel}
B.~Woodworth, S.~Gunasekar, J.~Lee, E.~Moroshko, P.~Savarese, I.~Golan, D.~Soudry, and N.~Srebro.
\newblock Kernel and rich regimes in overparameterized models.
\newblock In \emph{Conference on Learning Theory (COLT)}, 2020.

\bibitem[Geiger et~al.(2020)]{geiger2020disentangling}
M.~Geiger, A.~Jacot, S.~Spigler, F.~Gabriel, L.~Sagun, S.~d'Ascoli, G.~Biroli, C.~Hongler, and M.~Wyart.
\newblock Disentangling feature and lazy training in deep neural networks.
\newblock \emph{Journal of Statistical Mechanics: Theory and Experiment}, 2020(11):113301, 2020.

\bibitem[Saxe et~al.(2014)]{saxe2014exact}
A.~M.~Saxe, J.~L.~McClelland, and S.~Ganguli.
\newblock Exact solutions to the nonlinear dynamics of learning in deep linear neural networks.
\newblock In \emph{International Conference on Learning Representations (ICLR)}, 2014.

\bibitem[Du et~al.(2019)]{du2019gradient}
S.~S.~Du, X.~Zhai, B.~P{\'o}czos, and A.~Singh.
\newblock Gradient descent provably optimizes over-parameterized neural networks.
\newblock In \emph{International Conference on Learning Representations (ICLR)}, 2019.

\bibitem[Nguyen(2021)]{nguyen2021tight}
Q.~Nguyen.
\newblock On the proof of global convergence of gradient descent for deep {ReLU} networks with linear widths.
\newblock In \emph{International Conference on Machine Learning (ICML)}, 2021.

\bibitem[Cho \& Saul(2009)]{cho2009kernel}
Y.~Cho and L.~K.~Saul.
\newblock Kernel methods for deep learning.
\newblock In \emph{Advances in Neural Information Processing Systems (NeurIPS)}, 2009.

\bibitem[Novak et~al.(2019)]{novak2019neural}
R.~Novak, L.~Xiao, Y.~Bahri, J.~Lee, G.~Yang, J.~Hron, D.~A.~Abolafia, J.~Pennington, and J.~Sohl-Dickstein.
\newblock Bayesian deep convolutional networks with many channels are {G}aussian processes.
\newblock In \emph{International Conference on Learning Representations (ICLR)}, 2019.

\bibitem[Arora et~al.(2019)]{arora2019exact}
S.~Arora, S.~S.~Du, W.~Hu, Z.~Li, R.~Salakhutdinov, and R.~Wang.
\newblock On exact computation with an infinitely wide neural net.
\newblock In \emph{Advances in Neural Information Processing Systems (NeurIPS)}, 2019.

\bibitem[Fisher(1972)]{fisher1972scaling}
M.~E.~Fisher.
\newblock Scaling, universality and renormalization group theory.
\newblock In F.~J.~W.~Hahne, editor, \emph{Critical Phenomena}, Lecture Notes in Physics, pages 1--139. Springer, 1972.

\bibitem[Binder(1981)]{binder1981finite}
K.~Binder.
\newblock Finite size scaling analysis of {I}sing model block distribution functions.
\newblock \emph{Zeitschrift f{\"u}r Physik B Condensed Matter}, 43(2):119--140, 1981.

\bibitem[Keller(1977)]{keller1977numerical}
H.~B.~Keller.
\newblock Numerical solution of bifurcation and nonlinear eigenvalue problems.
\newblock In P.~H.~Rabinowitz, editor, \emph{Applications of Bifurcation Theory}, pages 359--384. Academic Press, 1977.

\bibitem[Allgower \& Georg(2003)]{allgower2003introduction}
E.~L.~Allgower and K.~Georg.
\newblock \emph{Introduction to Numerical Continuation Methods}.
\newblock SIAM, 2003.

\bibitem[Benettin et~al.(1980)]{benettin1980lyapunov}
G.~Benettin, L.~Galgani, A.~Giorgilli, and J.-M.~Strelcyn.
\newblock Lyapunov characteristic exponents for smooth dynamical systems and for {H}amiltonian systems; a method for computing all of them.
\newblock \emph{Meccanica}, 15(1):9--20, 1980.

\bibitem[Williams \& Seeger(2001)]{williams2001nystrom}
C.~K.~I.~Williams and M.~Seeger.
\newblock Using the {N}ystr{\"o}m method to speed up kernel machines.
\newblock In \emph{Advances in Neural Information Processing Systems (NeurIPS)}, 2001.

\bibitem[Gittens \& Mahoney(2016)]{gittens2016revisiting}
A.~Gittens and M.~W.~Mahoney.
\newblock Revisiting the {N}ystr{\"o}m method for improved large-scale machine learning.
\newblock \emph{Journal of Machine Learning Research}, 17(117):1--65, 2016.

\bibitem[Crandall \& Rabinowitz(1971)]{crandall1971bifurcation}
M.~G.~Crandall and P.~H.~Rabinowitz.
\newblock Bifurcation from simple eigenvalues.
\newblock \emph{Journal of Functional Analysis}, 8(2):321--340, 1971.

\end{thebibliography}

\end{document}
